\documentclass[11pt]{article}
\pdfinfoomitdate=1
\pdftrailerid{}
\pdfsuppressptexinfo=15
\usepackage[T1]{fontenc}
\usepackage{lmodern}
\usepackage[margin=1in]{geometry}
\usepackage{amsmath,amssymb,amsthm,mathtools}
\usepackage{enumitem,microtype}
\usepackage{xcolor}
\usepackage[colorlinks=true,linkcolor=blue!40!black,citecolor=blue!40!black,urlcolor=blue!40!black]{hyperref}
\input{glyphtounicode}
\pdfgentounicode=1
\pdfglyphtounicode{vextendsingle}{23D0}
\pdfglyphtounicode{vextenddouble}{2016}
\pdfglyphtounicode{slashbig}{002F}
\pdfglyphtounicode{parenleftBig}{0028}
\pdfglyphtounicode{parenrightBig}{0029}
\pdfglyphtounicode{mu}{03BC}

\pdfglyphtounicode{parenleftbig}{0028}
\pdfglyphtounicode{parenrightbig}{0029}
\pdfglyphtounicode{bracketleftbig}{005B}
\pdfglyphtounicode{bracketrightbig}{005D}
\pdfglyphtounicode{parenlefttp}{239B}
\pdfglyphtounicode{parenleftbt}{239D}
\pdfglyphtounicode{parenrighttp}{239E}
\pdfglyphtounicode{parenrightbt}{23A0}

\pdfglyphtounicode{braceleftbig}{007B}
\pdfglyphtounicode{bracerightbig}{007D}
\pdfglyphtounicode{parenleftbigg}{0028}
\pdfglyphtounicode{parenrightbigg}{0029}
\pdfglyphtounicode{braceleftbigg}{007B}
\pdfglyphtounicode{bracerightbigg}{007D}
\pdfglyphtounicode{parenleftBigg}{0028}
\pdfglyphtounicode{parenrightBigg}{0029}
\pdfglyphtounicode{braceleftBigg}{007B}
\pdfglyphtounicode{bracerightBigg}{007D}

\pdfglyphtounicode{bracketleftbigg}{005B}
\pdfglyphtounicode{bracketrightbigg}{005D}
\pdfglyphtounicode{bracketleftBigg}{005B}
\pdfglyphtounicode{bracketrightBigg}{005D}
\pdfglyphtounicode{bracketleftBig}{005B}
\pdfglyphtounicode{bracketrightBig}{005D}
\pdfglyphtounicode{floorleftBigg}{230A}
\pdfglyphtounicode{floorrightBigg}{230B}
\pdfglyphtounicode{ceilingleftBigg}{2308}
\pdfglyphtounicode{ceilingrightBigg}{2309}

\pdfglyphtounicode{bracelefttp}{23A7}
\pdfglyphtounicode{braceleftmid}{23A8}
\pdfglyphtounicode{braceleftbt}{23A9}
\pdfglyphtounicode{braceex}{23AA}

\pdfglyphtounicode{summationtext}{2211}
\pdfglyphtounicode{summationdisplay}{2211}
\pdfglyphtounicode{productdisplay}{220F}
\pdfglyphtounicode{producttext}{220F}
\pdfglyphtounicode{logicalortext}{22C1}
\pdfglyphtounicode{logicalordisplay}{22C1}
\pdfglyphtounicode{uniontext}{22C3}
\pdfglyphtounicode{uniondisplay}{22C3}
\pdfglyphtounicode{integraldisplay}{222B}
\pdfglyphtounicode{radicalbig}{221A}
\pdfglyphtounicode{radicalBig}{221A}
\pdfglyphtounicode{radicalbigg}{221A}
\pdfglyphtounicode{radicalBigg}{221A}

\pdfglyphtounicode{hatwide}{02C6}
\pdfglyphtounicode{tildewide}{02DC}

\pdfglyphtounicode{mapsto}{007C}

\hypersetup{
  pdftitle={Vanishing of higher bounded Hochschild cohomology},
  pdfauthor={OpenAI},
  pdfsubject={The Kadison--Ringrose cohomology conjecture}
}
\setlength{\emergencystretch}{1.5em}
\numberwithin{equation}{section}
\newtheorem{theorem}{Theorem}[section]
\newtheorem{lemma}[theorem]{Lemma}
\newtheorem{proposition}[theorem]{Proposition}
\newtheorem{corollary}[theorem]{Corollary}
\theoremstyle{definition}
\newtheorem{definition}[theorem]{Definition}
\theoremstyle{remark}
\newtheorem{remark}[theorem]{Remark}
\newcommand{\C}{\mathbb C}
\newcommand{\R}{\mathbb R}
\newcommand{\N}{\mathbb N}
\newcommand{\Z}{\mathbb Z}
\newcommand{\U}{\mathcal U}
\newcommand{\E}{\mathbb E}
\newcommand{\id}{\operatorname{id}}
\DeclarePairedDelimiter{\norm}{\lVert}{\rVert}
\DeclareMathOperator{\supp}{supp}
\DeclareMathOperator{\tr}{tr}
\DeclareMathOperator{\spanop}{span}
\title{Vanishing of higher bounded Hochschild cohomology}
\author{OpenAI}
\date{September 23, 2026}
\begin{document}
\maketitle
\begin{abstract}
We prove that every bounded Hochschild cocycle of degree at least two
on a complex von Neumann algebra, with values in the algebra itself,
has a bounded primitive. Together with the established degree-one
inner-derivation theorem, this resolves the Kadison--Ringrose cohomology
conjecture positively.
\end{abstract}

\section{Introduction}\label{sec:introduction}

The continuous Hochschild cohomology of an operator algebra measures
the obstruction to solving its bounded cocycle equations. The theory
also connects these equations with stability of multiplication under
small perturbations~\cite[Chapter~7]{SinclairSmith1995}.
For a von Neumann algebra acting on itself by multiplication, the
degree-one equation is the derivation equation, and its vanishing is
the inner-derivation theorem of Kadison and
Sakai~\cite{Kadison1966,Sakai1966}.
The higher-degree theory was developed by Johnson, Kadison, and
Ringrose~\cite{KadisonRingrose1971,KadisonRingrose1971II,JKR1972}.
Throughout, boundedness means ordinary
operator-norm boundedness of a multilinear map. No complete-boundedness
condition is imposed.

For a complex von Neumann algebra \(M\), let \(C_b^k(M,M)\) be the
Banach space of bounded complex \(k\)-linear maps \(M^k\to M\),
with norm
$\norm{f}=\sup_{\norm{a_1},\ldots,\norm{a_k}\leq1}
\norm{f(a_1,\ldots,a_k)}$, and put \(C_b^0(M,M)=M\).
For \(k\geq1\), the differential is
\begin{align}
(df)(a_1,\ldots,a_{k+1})
={}&a_1f(a_2,\ldots,a_{k+1})\notag\\
&+\sum_{j=1}^k(-1)^j
 f(a_1,\ldots,a_ja_{j+1},\ldots,a_{k+1})\notag\\
&+(-1)^{k+1}f(a_1,\ldots,a_k)a_{k+1}.
\label{eq:differential}
\end{align}
In degree zero, \((db)(a)=ab-ba\). As usual, \(d^2=0\).
We use the notation
\[
H_b^k(M,M)=
\ker\bigl(d:C_b^k(M,M)\to C_b^{k+1}(M,M)\bigr)
\big/
\operatorname{im}\bigl(d:C_b^{k-1}(M,M)\to C_b^k(M,M)\bigr).
\]
The denominator is the actual image, without a closure.

The Kadison--Ringrose cohomology conjecture asks whether
\(H_b^k(M,M)=0\) for every complex von Neumann algebra \(M\) and
every positive integer \(k\). Sinclair and Smith trace this
self-coefficient formulation to 1967~\cite[p.~3376]{SinclairSmith1998}.
Our main result settles the higher degrees.

\begin{theorem}\label{thm:main}
Let \(M\) be a complex von Neumann algebra and let \(k\geq2\).
Every \(f\in C_b^k(M,M)\) satisfying \(df=0\) has the form
\(f=dg\) for some \(g\in C_b^{k-1}(M,M)\).
Consequently \(H_b^k(M,M)=0\).
\end{theorem}

Together with the inner-derivation theorem,
Theorem~\ref{thm:main} resolves the Kadison--Ringrose cohomology
conjecture positively.

\subsection{Prior work}

Kadison and Ringrose proved vanishing for type~I and hyperfinite von
Neumann algebras~\cite[Theorem~4.4]{KadisonRingrose1971}
\cite[Theorem~3.1]{KadisonRingrose1971II}. Their extended-cobounding
theorem also gives a primitive for an $M$-valued cocycle on a general
concrete von Neumann algebra $M\subset B(H)$ when the primitive may
take values in
$B(H)$~\cite[Theorem~2.4]{KadisonRingrose1971II}; keeping the values in
the original algebra is the essential self-coefficient problem here.
Johnson, Kadison, and Ringrose then identified continuous cohomology
with separately normal cohomology~\cite[Theorem~5.6]{JKR1972}.
The completely bounded methods of Christensen, Effros, and
Sinclair~\cite{CES1987} established vanishing whenever the type
$\mathrm{II}_1$ central summand absorbs the hyperfinite
$\mathrm{II}_1$ factor, including all algebras without a type
$\mathrm{II}_1$ summand; see the precise restatement
in~\cite[p.~636, equation~(1.1)]{CPSS2003}.

The remaining finite-algebra problem has several substantial positive
cases. Sinclair and Smith proved vanishing for algebras with separable
predual and a Cartan subalgebra~\cite{SinclairSmith1998Cartan};
Cameron removed the separability hypothesis for
type~$\mathrm{II}_1$ factors with a Cartan maximal abelian
subalgebra~\cite[Theorem~3.2]{Cameron2009}.
Christensen, Pop, Sinclair, and Smith proved vanishing for all
type~$\mathrm{II}_1$ factors with property~$\Gamma$, including
nonseparable ones~\cite[Theorem~7.2]{CPSS2003}.
Qian and Shen extended this direction to type~$\mathrm{II}_1$
algebras with separable predual and their generalized
property~$\Gamma$~\cite[Theorem~6.4]{QianShen2015}.
In degree two, Pop and Smith proved vanishing for the tensor product
of any two type~$\mathrm{II}_1$ von Neumann algebras, with no
separability assumption~\cite[Theorem~4.5]{PopSmith2010}.
Theorem~\ref{thm:main} removes the structural restrictions on the
finite algebra in every degree $k\geq2$.

The Cartan argument of Sinclair and Smith also provides a useful
methodological comparison. They show that complete boundedness in the
last cocycle variable suffices. They use a projection of Christensen
and Sinclair~\cite{ChristensenSinclair1995} from completely bounded
linear maps onto right module maps, in the form recalled
in~\cite[Theorem~3.4]{SinclairSmith1998}, to obtain a primitive
\cite[Section~4]{SinclairSmith1998}.
Thus the passage from averaging one variable to a module identity
already has a central role in the earlier theory.

\subsection{The new ingredients}

We obtain the needed module identity by averaging over a single
unitary random walk. Averaging one cocycle input first gives only a
harmonic map; the main new input is a Liouville theorem that makes
such a map commute with right multiplication.
Suppose that \(M\) is a type \(\mathrm{II}_1\) algebra with separable
predual and faithful normal tracial state \(\tau\), and write
\(\norm{x}_2=\tau(x^*x)^{1/2}\). Let $\U(M)$ denote its unitary group.
We construct a symmetric, countably supported probability measure
\(\mu\) on \(\U(M)\), with a \(2\)-dense set of atoms, such that
\begin{equation}\label{eq:intro-harmonic}
T(x)=\E_{U\sim\mu}\bigl[T(xU)U^*\bigr]
\quad(x\in M)
\qquad\Longrightarrow\qquad
T(x)=T(1)x.
\end{equation}
Here \(T:M\to M\) is any bounded linear map that is \(2\)-norm
continuous on each operator-norm ball. No positivity or
complete-boundedness assumption is imposed.
The measure is fixed before the harmonic map is chosen.

For a separately normal cocycle, Akemann's continuity theorem
\cite[Corollary~III.10]{Akemann1967} supplies the stated continuity in
the last input. A uniform continuity modulus survives the averaging
limit, even though that limit need not remain normal, and also passes
to the tracial ultrapowers used in the Liouville proof.

The proof of \eqref{eq:intro-harmonic} combines probabilistic and
algebraic arguments. The finite free configurations come from Popa's
construction of free-independent sequences~\cite{Popa1995}, in the
formulation of his later ultraproduct independence
theorem~\cite{Popa2014}. These give finite unitary configurations with
prescribed small word traces. A hierarchy of increasingly rare,
increasingly large sets of such unitaries gives a single random walk
whose long independent products approach free Haar joint moments
along a subsequence. Only moment convergence is required.

Second, we use the classical harmonic-martingale construction
\cite[p.~415]{Furstenberg1963} to obtain an \(L^2(M,\tau)\)-valued
martingale. A nonconstant limit would produce free Haar unitaries
\(g_1,g_2,\ldots\) in a tracial ultrapower and distinct \(a,b\in M\).
The coordinatewise lift of \(T\) to that ultrapower sends each
nonempty signed word \(w\) with pairwise distinct unsigned indices
to \(aw\) or \(bw\), according to the sign of its first letter;
here \(a,b\) are identified with their constant images.
We prove that these identities are impossible for a bounded
map with the stated continuity.
The argument averages free generators, discards repeated surviving
indices using Haagerup's free-group length estimate~\cite{Haagerup1979},
and passes to a unitary
polar part. Positive and negative powers of that unitary would then
carry different left multipliers. Uniformly bounded sine sums and a
spectral partition force the multipliers to agree.
This rigidity statement, proved in
Theorem~\ref{thm:rigidity}, is independent of the cocycle problem.

Finally, in this tracial type \(\mathrm{II}_1\) setting, for a
separately normal cocycle \(f\) we take a pointwise ultraweak cluster
point of the Ces\`aro averages in its last input.
The Liouville theorem identifies the resulting fixed map as a
right module map in that input. Averaging the original cocycle equation
then gives a primitive \(g\) with
\(\norm{g}\leq\norm{f}\).
Together with compactness, this uniform estimate removes the
separability assumption. We then partition the type
\(\mathrm{II}_1\) summand into central pieces admitting faithful
normal tracial states. Lemma~\ref{lem:central-homotopy} corrects
mixed inputs from distinct pieces, and uniform bounds permit their
assembly. The classical vanishing theorem handles the complementary
summand.

\subsection{Conventions and classical reductions}

All von Neumann algebras and maps are complex, and subalgebras are
unital unless otherwise stated. A finite algebra is called
type \(\mathrm{II}_1\) here when it has no nonzero type~I central
summand; its center need not be trivial.
A map \(f:M^k\to M\) is \emph{separately normal} if it is
ultraweakly continuous in each input with the others fixed.
For a finite algebra with faithful normal tracial state, the
\(2\)-norm on bounded sets determines the strong-star topology.
The inherited traces on subalgebras and tracial ultrapowers will
usually also be denoted by \(\tau\).

We use the following established cohomology results.
\begin{samepage}
\begin{proposition}[Classical reductions]\label{prop:classical}
Let \(M\) be a von Neumann algebra and \(k\geq2\).
\begin{enumerate}[label=\textup{(\roman*)}]
\item If \(f\in C_b^k(M,M)\) is a cocycle, there are a separately
normal bounded cocycle \(f_0\) and a bounded \((k-1)\)-cochain \(b\)
such that \(f=f_0+db\).
\item If \(M\) has no type \(\mathrm{II}_1\) central summand, then
\(H_b^k(M,M)=0\).
\end{enumerate}
\end{proposition}
\end{samepage}

Part~\textup{(i)} is the Johnson--Kadison--Ringrose normal-reduction
theorem~\cite[Lemma~5.4 and Theorem~5.6]{JKR1972}; its formulation for self-coefficients is recalled
in \cite[Section~2]{CPSS2003}.
Part~\textup{(ii)} combines the classical type~I and properly infinite
vanishing results; see \cite[Theorem~4.4]{KadisonRingrose1971},
\cite{CES1987}, and \cite[p.~636, equation~(1.1)]{CPSS2003}.
Both statements concern ordinary continuous cohomology and the actual
image of \(d\). These are the only cohomological inputs in the proof.

Section~\ref{sec:walk} constructs the walk.
Section~\ref{sec:liouville} proves its Liouville property using
Theorem~\ref{thm:rigidity}, proved in Section~\ref{sec:rigidity}.
Section~\ref{sec:cohomology} constructs bounded primitives and
completes the reductions to arbitrary \(M\).

\section{A random walk with asymptotically free products}\label{sec:walk}

Throughout this section, $(M,\tau)$ is a finite von Neumann algebra of
type $\mathrm{II}_1$ with separable predual and faithful normal tracial
state.  In particular, $M$ may have a nontrivial center.  We construct a
single symmetric probability measure on $\U(M)$ whose independent long
products approach the joint moments of free Haar unitaries along a
subsequence.  The construction combines finite approximation to freeness
with increasingly rare increments chosen at successive stages.

Recall that unital subalgebras of a tracial von Neumann algebra are
\emph{free} if every alternating product of centered elements from
successively different subalgebras has trace zero.  Here centered means
having trace zero.  A unitary $u$ is \emph{Haar} if
$\tau(u^k)=0$ for every $k\in\Z\setminus\{0\}$.  A free Haar family
$(u_i)$ is equivalently characterized by
\[
 \tau\bigl(w(u_1,\ldots,u_r)\bigr)=0
 \qquad(w\in\mathbb F_r\setminus\{e\}),
\]
for every finite subfamily, where $\mathbb F_r$ denotes the free group on
$r$ specified generators.  The equivalence follows by expanding centered
Laurent polynomials and then taking bounded $2$-norm approximations.
The length $|w|$ is its reduced word length.

For a probability measure $\mu$ on $\U(M)$, let $\mu^{*n}$ be the law
of the ordered product $U_1\cdots U_n$ of independent increments of
law $\mu$.  Symmetry means that $U^{-1}$ has the same law as $U$.

\begin{theorem}\label{thm:walk}
Let $M$ be a finite type $\mathrm{II}_1$ von Neumann algebra with
separable predual and faithful normal tracial state $\tau$.  There
exist a symmetric, countably supported probability measure $\mu$ on
$\U(M)$ and a strictly increasing sequence $n_\ell\to\infty$ with the
following properties.

The positive-mass atoms of $\mu$ contain a countable $2$-norm dense
subset of $\U(M)$, and hence generate $M$ as a von Neumann algebra.
For every $r\in\N$, every nonempty finite set
$\mathcal W\subset\mathbb F_r\setminus\{e\}$, and every $\eta>0$, if
$G_1^{(\ell)},\ldots,G_r^{(\ell)}$ are independent with law
$\mu^{*n_\ell}$, then
\[
 \Pr\left(
   \max_{w\in\mathcal W}
   \left|\tau\bigl(w(G_1^{(\ell)},\ldots,G_r^{(\ell)})\bigr)\right|
   \geq\eta
 \right)\longrightarrow0.
\]
Thus these products have joint free Haar moments in probability.
\end{theorem}

We first obtain finite unitary configurations for which prescribed
words in path products have small traces.  We then place such
configurations at increasingly rare levels of a single measure,
arranging that each finite moment test succeeds with probability
tending to one.

\subsection{Finite requirements from ultrapower freeness}

The freeness input originates in Popa's free-independent sequence
construction~\cite[Lemma~1.4 and Theorem~2.1]{Popa1995}.
We use the following special case of its later ultraproduct
formulation~\cite[Theorem~0.1(a)]{Popa2014}.

For a tracial von Neumann algebra $(P,\tau_P)$ and a free ultrafilter
$\omega$ on $\N$, its tracial ultrapower is
\begin{equation}\label{eq:tracial-ultrapower}
 P^\omega=\ell^\infty(\N,P)/\mathcal I_\omega,
 \qquad
 \mathcal I_\omega
 =\{(x_m):\lim_{m\to\omega}\norm{x_m}_2=0\}.
\end{equation}
It is a finite von Neumann algebra with faithful normal trace
$\tau_P^\omega((x_m)^\omega)=\lim_{m\to\omega}\tau_P(x_m)$.
We identify $P$ with its constant sequences.

\begin{lemma}[Popa's ultrapower freeness]\label{lem:popa-special}
Let $P$ be a $\mathrm{II}_1$ factor and let $B\subset P^\omega$ be a
unital von Neumann subalgebra with separable predual.  There is a diffuse
abelian von Neumann subalgebra of $P^\omega$ free from $B$.
Consequently, for every $d\in\N$, there are Haar unitaries
$u_1,\ldots,u_d\in P^\omega$ such that
$B,W^*(u_1),\ldots,W^*(u_d)$ are free.
\end{lemma}

\begin{proof}
Apply Theorem~0.1(a) of \cite{Popa2014} with every coordinate algebra
and every coordinate subalgebra equal to $P$.  The required
non-intertwining condition holds: a nonzero corner of the diffuse factor
$P$ cannot embed in the scalar relative commutant, or in a finite matrix
amplification of it.  Moreover, $P^\omega$ is a factor, so the relative
commutant of the full ultraproduct subalgebra is $\C1$.  The conclusion
for the separable centered subspace $B\ominus\C1$ is therefore ordinary
scalar-trace freeness. This is the case in which every coordinate
subalgebra equals its ambient factor, discussed after that theorem
and in \cite[Sections~1.5--1.6]{Popa2014}.

A diffuse abelian tracial algebra contains a Haar unitary.  Choose one,
adjoin it to $B$, and repeat.  Finite adjunction preserves separable
predual: a countable generating set gives a separable tracial Hilbert
space by polynomial approximation.  Thus all $d$ choices take place in
the same ultrapower.  At each step the new algebra is free from the
algebra generated at the preceding steps, giving the asserted family.
\end{proof}

We next record why a single fresh label in each path is sufficient.
A \emph{path expression} is a finite ordered product of fixed unitary
coefficients and formal letters $U_j^{\pm1}$.  Occurrences and labels in
such expressions are counted before any simplification.

\begin{lemma}[Distinguished letters]\label{lem:distinguished-letters}
Let $B,W^*(u_1),\ldots,W^*(u_d)$ be free subalgebras of a tracial von
Neumann algebra, with each $u_j$ Haar.  Let
$y_1,\ldots,y_r$ be path expressions evaluated at these unitaries, with
coefficients in $\U(B)$.  Suppose that for each $i$ there is a label
$h_i$ occurring exactly once in the entire list of path expressions,
and that its occurrence belongs to $y_i$.  Then
\[
 \tau\bigl(w(y_1,\ldots,y_r)\bigr)=0
 \qquad(w\in\mathbb F_r\setminus\{e\}).
\]
\end{lemma}

\begin{proof}
The labels $h_1,\ldots,h_r$ are distinct.  Put all nondistinguished
unitaries into the coefficient algebra:
\[
 B_0=W^*\bigl(B,u_j:j\notin\{h_1,\ldots,h_r\}\bigr),
 \qquad
 C=W^*(u_{h_1},\ldots,u_{h_r}).
\]
Then $B_0$ and $C$ are free.  This grouping property follows directly
by expanding centered polynomials in a free family into alternating
centered products; bounded approximation extends it to the generated
von Neumann algebras.  Write
\[
 y_i=A_i u_{h_i}^{\epsilon_i}B_i,
 \qquad A_i,B_i\in\U(B_0),\quad \epsilon_i\in\{-1,1\}.
\]
If $w=x_{i_1}^{\sigma_1}\cdots x_{i_t}^{\sigma_t}$ is nonempty and
reduced, its evaluated product has the form
\[
 c_0 u_{h_{i_1}}^{\epsilon_{i_1}\sigma_1}c_1
 \cdots u_{h_{i_t}}^{\epsilon_{i_t}\sigma_t}c_t,
 \qquad c_0,\ldots,c_t\in\U(B_0).
\]
The displayed distinguished-letter string is itself reduced.  Indeed,
an adjacent inverse pair would require $i_j=i_{j+1}$ and
$\sigma_j=-\sigma_{j+1}$, contrary to the reduction of $w$.

Expand every $c_j$, including the two outside coefficients, as its
scalar part plus its centered part.  In each resulting term, removing
scalar coefficients joins distinguished letters into nonempty
contiguous subwords of that reduced string.  Every such block is a
centered element of $C$.  The remaining factors alternate between
centered elements of $B_0$ and centered elements of $C$, and there is
at least one $C$ factor.  Freeness makes the trace zero.  If there are
no centered coefficient factors, the term is just a scalar multiple of
a nonempty reduced Haar word and again has trace zero.  This proves
the assertion, with no condition on the first and last indices of $w$.
\end{proof}

\begin{lemma}[Simultaneous finite path requirements]\label{lem:finite-path-tests}
Fix a finite set $F\subset\U(M)$, an integer $d\geq1$, and
$\varepsilon>0$.  Consider a finite collection of tests.  Each test
consists of a finite list of path expressions $y_1,\ldots,y_r$ in the
alphabet $F\cup\{U_j^{\pm1}:1\leq j\leq d\}$ and a nonidentity word
$w\in\mathbb F_r$.  Assume that each path in each test has a
distinguished label occurring exactly once in that test's entire list,
as in Lemma~\ref{lem:distinguished-letters}.  The distinguished labels
may differ between tests.  There are $u_1,\ldots,u_d\in\U(M)$ for which
all the tests satisfy
\[
 \left|\tau\bigl(w(y_1(u),\ldots,y_r(u))\bigr)\right|
 <\varepsilon.
\]
\end{lemma}

\begin{proof}
First suppose that $M=P$ is a factor.  The constant coefficient algebra
$W^*(F)\subset P^\omega$ has separable predual.  Apply
Lemma~\ref{lem:popa-special} and then
Lemma~\ref{lem:distinguished-letters}.  The same free $d$-tuple makes
every prescribed trace exactly zero in $P^\omega$, even though different
tests can designate different labels.

Every unitary of $P^\omega$ has unitary representatives.  To see this,
take a bounded self-adjoint logarithm there, lift it to a bounded
sequence of self-adjoint elements of $P$, and exponentiate coordinatewise.
Choose such representatives for the $d$ unitaries.  Each test trace
converges to zero along $\omega$.  Since the family of tests is finite,
the coordinate sets on which their absolute values are less than
$\varepsilon$ have a common member.  Its $d$ coordinate unitaries solve
all the requirements in $P$.

For general $M$, the factor case will give successful tuples in
almost every fiber.  The remaining task is to choose them measurably.
Use the central direct integral
\cite[Sections~III.1.6 and III.5.1]{Blackadar2006}
\[
 M=\int_\Omega^\oplus M_s\,d\nu(s),
 \qquad
 \tau(x)=\int_\Omega\tau_s(x(s))\,d\nu(s),
\]
where $(\Omega,\nu)$ is a probability space and, almost everywhere,
$M_s$ is a $\mathrm{II}_1$ factor with its normalized trace $\tau_s$.
Standard separable direct-integral theory provides countably many
bounded measurable sections $a_j$ whose values generate $M_s$ for
almost every $s$.  After removing one null set, all these properties
and the unitary identities for the finitely many coefficients hold
simultaneously.

There is a countable family $(v_j)$ of measurable unitary sections
which is dense in each fiber unitary group for
$\norm{\cdot}_{2,s}$.  Here is an explicit construction.  Enumerate the
sections
\[
 \exp(iq(a_1,a_2,\ldots)),
\]
where $q$ ranges over self-adjoint $*$-polynomials in finitely many
variables with coefficients in $\mathbb Q+i\mathbb Q$.  This includes
symmetrizations of arbitrary rational $*$-polynomials.  In a fixed
fiber, any unitary is $\exp(ih)$ for a bounded self-adjoint $h$.
Kaplansky density and norm approximation by rational polynomials
approximate $h$ in $2$-norm by the indicated self-adjoint polynomials.
For bounded self-adjoint $a,b$, the integral identity
\[
 e^{ia}-e^{ib}
 =i\int_0^1 e^{i(1-t)a}(a-b)e^{itb}\,dt
\]
implies $\norm{e^{ia}-e^{ib}}_2\leq\norm{a-b}_2$.
Thus these sections have the required fiberwise density.  Each section
is globally bounded, and exponentiation makes its norm one.

Enumerate all ordered $d$-tuples drawn from $(v_j)$ as
$\boldsymbol v^{(1)},\boldsymbol v^{(2)},\ldots$.  Let $E_k$ be the
set of fibers where candidate $k$ satisfies every test with absolute
trace less than $\varepsilon/2$.  These sets are measurable, since
the traces of products of measurable sections are measurable.  They
cover almost every fiber.  Indeed, the factor case gives a successful
tuple in each such fiber, and the requirements are open in the product
$2$-norm topology.  For this last assertion, telescoping products of
unitaries gives
\[
 \norm{z_1\cdots z_L-z'_1\cdots z'_L}_2
 \leq\sum_{j=1}^L\norm{z_j-z'_j}_2,
 \qquad |\tau_s(x)|\leq\norm{x}_{2,s},
\]
and inversion is also a $2$-norm isometry on unitaries.  The countable
dense tuples therefore meet the open set of successful choices.

Take the first successful tuple: on the disjoint measurable sets
\[
 E_k\setminus\bigcup_{j<k}E_j
\]
use $\boldsymbol v^{(k)}$, and use the identity tuple on the null
complement.  The resulting sections $u_1,\ldots,u_d$ are measurable
and unitary, so belong to $M$.  Each test has fiber trace of absolute
value less than $\varepsilon/2$ almost everywhere.  Integrating gives
an absolute global trace at most $\varepsilon/2<\varepsilon$, as
required.
\end{proof}

\subsection{One measure satisfying all moment tests}

\begin{proof}[Proof of Theorem~\ref{thm:walk}]
Choose a $2$-norm dense sequence $(v_j)_{j\geq1}$ in $\U(M)$ and put
\[
 \mu_0=\sum_{j\geq1}2^{-j-1}(\delta_{v_j}+\delta_{v_j^{-1}}).
\]
For $\ell\geq1$, fix the numerical parameters
\begin{align*}
 p_\ell&=2^{-2^\ell},
 &p_0&=1-\sum_{\ell\geq1}p_\ell,\\
 n_\ell&=p_\ell^{-3/2}=2^{3\cdot2^{\ell-1}},
 &K_\ell&=\ell n_\ell,\\
 d_\ell&=K_\ell^4.
\end{align*}
Here $n_\ell$ and $d_\ell$ are positive integers, and $p_0>0$ because
$p_\ell\leq4^{-\ell}$.  At time $n_\ell$ we will sample $\ell$ paths,
choosing level $j$ with probability $p_j$ at each increment.  These
parameters make level $\ell$ appear in every path with high probability
while later levels remain unlikely.  The much larger label count
$d_\ell$ makes repeated level-$\ell$ labels unlikely across the
$K_\ell$ increments.  We construct finite symmetric measures
\[
 \mu_\ell=\frac1{2d_\ell}
   \sum_{j=1}^{d_\ell}(\delta_{u_{\ell,j}}
                            +\delta_{u_{\ell,j}^{-1}})
\]
successively, and finally set
\[
 \mu=p_0\mu_0+\sum_{\ell\geq1}p_\ell\mu_\ell.
\]

At stage $\ell$, choose a finite symmetric set $F_\ell$ of base atoms
with $\mu_0(F_\ell)>1-K_\ell^{-2}$.  The coefficient alphabet consists
of $F_\ell$ and all $u_{j,h}^{\pm1}$ from stages $j<\ell$; it is
finite.  Introduce $d_\ell$ unknown unitary labels.  Consider every
ordered list of $\ell$ paths of length $n_\ell$ in this alphabet and
the signed unknown labels such that:
\begin{enumerate}
 \item each path contains at least one unknown label;
 \item no unknown label occurs twice anywhere in the entire list,
       irrespective of sign.
\end{enumerate}
For each such list and every nonempty reduced word of length at most
$\ell$ on $\ell$ generators, require its evaluated trace to have
absolute value less than $1/\ell$.  There are finitely many lists and
finitely many words: all alphabet sizes, path lengths, and word lengths
have been fixed before making this choice.  Designating any one unknown
occurrence in each path puts every test under
Lemma~\ref{lem:finite-path-tests}.  Apply that lemma to choose
$u_{\ell,1},\ldots,u_{\ell,d_\ell}$ satisfying all these requirements
simultaneously.  This completes the induction.

We estimate probabilities on a labeled sampling space.  An increment
first chooses a level $j$ with probability $p_j$.  At level zero it
chooses an atom with law $\mu_0$; at level $j\geq1$ it chooses one of
$d_j$ labels uniformly and an independent uniform sign.  Its value is
then the corresponding unitary.  This description gives law $\mu$ even
when distinct labels happen to represent the same element of $M$.

Sample $\ell$ independent paths of length $n_\ell$, involving
$K_\ell$ increments in all.  Let $\mathcal B_\ell$ be the union of
the following four events: a level greater than $\ell$ appears; a
level-zero increment lies outside $F_\ell$; a level-$\ell$ label is
repeated; or some path has no level-$\ell$ increment.  Union bounds give
\begin{align}
 \Pr(\mathcal B_\ell)
 &\leq K_\ell\sum_{j>\ell}p_j
       +K_\ell p_0\mu_0(\U(M)\setminus F_\ell)\notag\\
 &\quad+\frac{\binom{K_\ell}{2}}{d_\ell}
       +\ell(1-p_\ell)^{n_\ell}.
 \label{eq:walk-bad-events}
\end{align}
For the collision term, any specified pair of increments has
probability $p_\ell^2/d_\ell\leq1/d_\ell$ of using the same
level-$\ell$ label; the signs do not matter.  For the last term, a
specified path misses that level with probability
$(1-p_\ell)^{n_\ell}$.

Every term on the right of \eqref{eq:walk-bad-events} tends to zero.
Indeed, $p_{\ell+1}=p_\ell^2$ and $p_\ell\leq1/4$, so
\[
 \sum_{j>\ell}p_j
 \leq\frac{p_\ell^2}{1-p_\ell^2}
 \leq\frac{16}{15}p_\ell^2.
\]
Consequently, the four terms admit the respective bounds
\begin{align*}
 \frac{16}{15}\ell\sqrt{p_\ell},\qquad
 K_\ell^{-1},\qquad
 \frac{1}{2K_\ell^2},\qquad
 \ell\exp(-p_\ell^{-1/2}),
\end{align*}
all converging to zero.

Outside $\mathcal B_\ell$, the sampled paths form exactly one of the
lists tested at stage $\ell$.  All their earlier-level factors belong
to the coefficient alphabet, each path has a level-$\ell$ label, and
all occurrences of those labels are distinct.  Thus, simultaneously
for every $w\in\mathbb F_\ell$ with $1\leq|w|\leq\ell$, the trace
of the evaluated word has absolute value less than $1/\ell$.

Now fix $r$, a nonempty finite $\mathcal W$, and $\eta>0$.  For
$\ell\geq\max\{r,\max_{w\in\mathcal W}|w|\}$, append independent
paths to the given $r$ paths to obtain $\ell$ paths.  On the preceding
good event every word in $\mathcal W$ has absolute trace below
$1/\ell$.  Once $1/\ell<\eta$, the probability in the theorem is at
most $\Pr(\mathcal B_\ell)$ and therefore tends to zero.

The measure $\mu$ is symmetric and countably supported by construction.
Its base component has positive weight, so each $v_j$ is a
positive-mass atom.  Since bounded $2$-norm convergence implies strong
convergence in the tracial representation, the von Neumann algebra
generated by these atoms contains every unitary and hence is $M$.
Finally, the identity group word always has trace one, and every
$*$-monomial in unitaries reduces to a group word.  The asserted
convergence therefore describes all joint $*$-moments.
\end{proof}

\begin{remark}\label{rem:walk-only-moments}
Theorem~\ref{thm:walk} asserts convergence of scalar moments.  It does
not assert convergence of operator norms of polynomials in the sampled
products.  Its finite-family formulation permits further conditions
whose failure probabilities tend to zero to be imposed simultaneously
at any fixed stage of a later selection argument.
\end{remark}

\section{A Liouville theorem for the walk}\label{sec:liouville}

We show that a bounded harmonic map for the walk of
Theorem~\ref{thm:walk}, subject to continuity on bounded sets in the tracial
$2$-norm, is left multiplication by its value at $1$.  The probabilistic
part of the proof produces free Haar unitaries such that, on words with
pairwise distinct unsigned indices, a single map acts by one of two left
multipliers according to the sign of the first letter.
Theorem~\ref{thm:rigidity}, proved in Section~\ref{sec:rigidity},
forces those multipliers to coincide for bounded linear maps with the
same continuity in any finite von Neumann algebra with a faithful
normal tracial state.

\subsection{Continuity and tracial ultrapowers}

\begin{definition}\label{def:ball-continuity}
Let $(N,\rho)$ be a finite von Neumann algebra with a faithful normal
tracial state.  A linear map $R:N\to N$ is \emph{bounded-ball
$2$-continuous} if
\begin{equation}\label{eq:ball-sequential}
 \sup_j\norm{x_j}<\infty,\qquad \norm{x_j}_2\longrightarrow0
 \quad\Longrightarrow\quad
 \norm{R(x_j)}_2\longrightarrow0.
\end{equation}
Here and below $\norm{x}_2=\rho(x^*x)^{1/2}$ when the ambient algebra is
$(N,\rho)$.
\end{definition}

The following lemma records the continuity needed after each of the two
ultrapowers in the argument.  It applies without a separability
assumption.

\begin{lemma}\label{lem:ultrapower}
Let $(N,\rho)$ be as in Definition~\ref{def:ball-continuity}, and let
$R:N\to N$ be bounded and linear.  Condition~\eqref{eq:ball-sequential}
is equivalent to the following uniform assertion: for every $r<\infty$
and $\varepsilon>0$ there is $\delta>0$ such that
\begin{equation}\label{eq:ball-modulus}
 \norm{x}\le r,\quad \norm{x}_2<\delta
 \quad\Longrightarrow\quad \norm{R(x)}_2<\varepsilon.
\end{equation}
If this condition holds and $\omega$ is a free ultrafilter on $\N$, then
\begin{equation}\label{eq:ultrapower-map}
 R^\omega:N^\omega\longrightarrow N^\omega,
 \qquad R^\omega([(x_j)]_\omega)=[(R(x_j))]_\omega,
\end{equation}
is a well-defined bounded linear map with
$\norm{R^\omega}\le\norm{R}$.  It is bounded-ball $2$-continuous for
the ultrapower trace.  In particular, for every free ultrafilter $\nu$,
the map $(R^\omega)^\nu$ on $(N^\omega)^\nu$ has the same properties.
\end{lemma}

\begin{proof}
If~\eqref{eq:ball-modulus} fails for some $r,\varepsilon$, choose
$x_j$ with $\norm{x_j}\le r$, $\norm{x_j}_2<1/j$, and
$\norm{R(x_j)}_2\ge\varepsilon$.  This contradicts
\eqref{eq:ball-sequential}; the converse is immediate.  Applying the
uniform assertion to differences gives uniform $2$-norm continuity on
each fixed operator norm ball.

Use the tracial ultrapower and its ideal $\mathcal I_\omega$ from
\eqref{eq:tracial-ultrapower}, now for $(N,\rho)$.  Its tracial norm is
\[
 \norm{[(x_j)]_\omega}_2=\lim_{j\to\omega}\norm{x_j}_2.
\]
By~\eqref{eq:ball-modulus}, the coordinate map induced by $R$ preserves
$\mathcal I_\omega$.  Applying this observation to the difference of two
bounded representatives proves that~\eqref{eq:ultrapower-map} is
well-defined.  The quotient norm gives
\[
 \norm{R^\omega(X)}
 \le \norm{R}\inf_{[(x_j)]_\omega=X}\sup_j\norm{x_j}
 =\norm{R}\norm{X}.
\]

To transfer the modulus, fix $r<\infty$ and $\varepsilon>0$, and choose
$\delta$ in~\eqref{eq:ball-modulus} for the radius $r+1$ and the output
tolerance $\varepsilon/2$.  If $\norm{X}\le r$, the quotient norm
allows a representative with $\sup_j\norm{x_j}\le r+1$.  When
$\norm{X}_2<\delta/2$, the set
\[
 \{j:\norm{x_j}_2<\delta\}
\]
belongs to $\omega$.  Consequently
\[
 \norm{R^\omega(X)}_2
 =\lim_{j\to\omega}\norm{R(x_j)}_2
 \le\varepsilon/2<\varepsilon.
\]
Thus $\delta/2$ is a valid input tolerance for $R^\omega$ on the
radius-$r$ ball.  The proof used no separability of $N$ and applies again
with $N^\omega$ in its place, proving the last assertion.
\end{proof}

\subsection{The harmonic map and its martingale}

The probabilistic step uses the classical passage from harmonic
functions to martingales along random products; see
\cite[p.~415]{Furstenberg1963}. Here the function is
$D(g)=T(g)g^*$ with values in $L^2(M,\tau)$, and the required
convergence follows directly from orthogonality of Hilbert-space
martingale increments.

\begin{theorem}\label{thm:liouville}
Let $(M,\tau)$ be a type $\mathrm{II}_1$ von Neumann algebra with
separable predual and a faithful normal tracial state, and let $\mu$ be
the measure supplied by Theorem~\ref{thm:walk}.  Suppose that
$T:M\to M$ is bounded, linear, and bounded-ball $2$-continuous, and that
\begin{equation}\label{eq:harmonic-map}
 T(x)=\E_{U\sim\mu}\bigl[T(xU)U^*\bigr]
 \qquad(x\in M).
\end{equation}
Then
\begin{equation}\label{eq:liouville-conclusion}
 T(x)=T(1)x\qquad(x\in M).
\end{equation}
\end{theorem}

\begin{proof}
Write $C=\norm{T}$ and $H=L^2(M,\tau)$.  Expectations in
\eqref{eq:harmonic-map} are norm-convergent sums, since $\mu$ is
countably supported and $\norm{T(xU)U^*}\le C\norm{x}$.  For
$g\in\U(M)$ put
\[
 D(g)=T(g)g^*.
\]
Let $U_1,U_2,\ldots$ be independent with law $\mu$, and set
$W_0=1$, $W_n=U_1\cdots U_n$, and
$\mathcal F_n=\sigma(U_1,\ldots,U_n)$.  The $H$-valued random
variables $Z_n=D(W_n)$ have countable range and satisfy
$\norm{Z_n}_2\le\norm{Z_n}\le C$.  Moreover,
\begin{align*}
 \E[Z_{n+1}\mid\mathcal F_n]
 &=\E\bigl[T(W_nU_{n+1})U_{n+1}^*W_n^*
                  \mid\mathcal F_n\bigr] \\
 &=T(W_n)W_n^*=Z_n.
\end{align*}
Thus $(Z_n)$ is a martingale in $L^2(\Omega;H)$.  Its increments are
orthogonal, and
\[
 \norm{Z_0}_2^2+
 \sum_{j=0}^{n-1}\E\norm{Z_{j+1}-Z_j}_2^2
 =\E\norm{Z_n}_2^2\le C^2.
\]
It follows that $Z_n$ converges in $L^2(\Omega;H)$ to a random variable
$Z_\infty$, and continuity of conditional expectation gives
$Z_n=\E[Z_\infty\mid\mathcal F_n]$.  In particular, with
\begin{equation}\label{eq:martingale-error}
 e_k=\E\norm{Z_\infty-Z_k}_2^2\longrightarrow0,
\end{equation}
orthogonality gives the uniform Cauchy estimate
\begin{equation}\label{eq:martingale-cauchy}
 \E\norm{Z_n-Z_k}_2^2=e_k-e_n\le e_k
 \qquad(n\ge k).
\end{equation}

The operator norm ball
$K_C=\{x\in M:\norm{x}\le C\}$ is closed in $H$.  Indeed, from a
$2$-norm-convergent sequence in $K_C$, ultraweak compactness provides
an ultraweakly convergent subnet with limit in $K_C$.  Testing against
$\tau(y^*\,\cdot)$, for $y\in M$, identifies this limit with the
$H$-limit, since $M$ is dense in $H$.  A subsequence of $Z_n$ converges
almost surely in $H$, so $Z_\infty\in K_C$ almost surely.

If $Z_\infty=c$ almost surely, conditional expectation gives
$Z_0=Z_1=c$ almost surely.  Hence $c=T(1)$ and $D(u)=c$ at every atom
$u$ of $\mu$.  These atoms contain a $2$-norm-dense subset of
$\U(M)$ by Theorem~\ref{thm:walk}.  For any $u\in\U(M)$, choose
atoms $u_j$ with $\norm{u_j-u}_2\to0$.  Bounded-ball continuity gives
$T(u_j)\to T(u)$ in $2$-norm, while
$\norm{cu_j-cu}_2\le\norm{c}\norm{u_j-u}_2\to0$.  Since
$T(u_j)=cu_j$, we obtain $T(u)=cu$.  Every element of a unital complex
$C^*$-algebra is a finite linear combination of unitaries, proving
\eqref{eq:liouville-conclusion} in this case.

Suppose, for a contradiction, that $Z_\infty$ is not deterministic.
Its law is a Borel probability measure on the separable Hilbert space
$H$, so its support contains two distinct points $a,b$.  Because
$K_C$ is closed and has full measure, $a,b\in K_C\subset M$.

\medskip\noindent
\emph{Forward and reverse endpoints.}
We claim that, for every $\delta>0$, there are numbers
$\eta_\delta>0$ and $N_\delta$ such that, whenever $n\ge N_\delta$
and $G$ has law $\mu^{*n}$,
\begin{equation}\label{eq:two-endpoints}
 \Pr\bigl(\norm{D(G)-a}_2<\delta,
          \ \norm{D(G^{-1})-b}_2<\delta\bigr)
 \ge\eta_\delta.
\end{equation}
Realize $G=U_1\cdots U_n$, let $h=\lfloor n/2\rfloor$, and put
\[
 A_n=D(U_1\cdots U_h),\qquad
 B_n=D(U_n^{-1}\cdots U_{n-h+1}^{-1}).
\]
The two variables use disjoint sets of increments and are independent.
Symmetry of $\mu$ implies that both have the law of $Z_h$.  Apply
\eqref{eq:martingale-cauchy} to the forward string and separately to
the reversed inverse string, which is also a string of independent
increments with law $\mu$.  It gives
\begin{equation}\label{eq:endpoint-replacement}
 \E\norm{D(G)-A_n}_2^2\le e_h,
 \qquad
 \E\norm{D(G^{-1})-B_n}_2^2\le e_h.
\end{equation}
For fixed $\delta$, the support property gives
\[
 p_a=\Pr(\norm{Z_\infty-a}_2<\delta/3)>0,
 \qquad
 p_b=\Pr(\norm{Z_\infty-b}_2<\delta/3)>0.
\]
Mean-square convergence implies, for all sufficiently large $h$,
\[
 \Pr(\norm{Z_h-a}_2<2\delta/3)\ge p_a/2,
 \qquad
 \Pr(\norm{Z_h-b}_2<2\delta/3)\ge p_b/2.
\]
Independence of $A_n,B_n$, the triangle inequality, and Markov's
inequality therefore bound the probability in
\eqref{eq:two-endpoints} from below by
\[
 \frac{p_ap_b}{4}-\frac{18e_h}{\delta^2}.
\]
For large $n$ this is at least $p_ap_b/8$, proving the claim.  The
independence used here is that of the disjoint prefixes $A_n,B_n$;
the variables $G$ and $G^{-1}$ remain dependent.

\medskip\noindent
\emph{Concatenation of independent blocks.}
Let $G_1,\ldots,G_m$ be independent with law $\mu^{*n}$, and consider
\[
 v(G)=G_{i_1}^{\sigma_1}\cdots G_{i_t}^{\sigma_t},
 \qquad \sigma_j\in\{1,-1\},
\]
where $i_1,\ldots,i_t$ are pairwise distinct.  Each inverse block can
be represented by reversing and inverting its own increments.  By
symmetry, concatenating these distinct signed blocks is a length-$tn$
walk with independent increments of law $\mu$.  In particular,
\[
 \bigl(D(v(G)),D(G_{i_1}^{\sigma_1})\bigr)
 \quad\hbox{has the law of}\quad (Z_{tn},Z_n).
\]
Consequently
\begin{equation}\label{eq:independent-blocks}
 \E\norm{D(v(G))-D(G_{i_1}^{\sigma_1})}_2^2\le e_n.
\end{equation}
The restriction is on unsigned indices: using both $G_i$ and $G_i^{-1}$
would reuse a block and would not justify this estimate.

\medskip\noindent
\emph{Finite requirements and diagonal selection.}
Fix an integer $m\ge1$ first, and use tolerance $1/m$ in
\eqref{eq:two-endpoints}.  For $m$ independent length-$n$ products,
let $E_{m,n}$ be the event
\[
 \norm{D(G_i)-a}_2<1/m,\qquad
 \norm{D(G_i^{-1})-b}_2<1/m
 \quad(1\le i\le m).
\]
Independence between the paths gives numbers $c_m>0$ and $N_m$ such
that
\begin{equation}\label{eq:conditioning-mass}
 \Pr(E_{m,n})\ge c_m\qquad(n\ge N_m).
\end{equation}
For example, one may take $c_m=(\eta_{1/m})^m$.

Let $\mathcal Q_m$ be the finite set of nonidentity reduced words of
length at most $m$ on the symbols $x_1,\ldots,x_m$.  Let
$\mathcal V_m$ be the finite set of all nonempty signed words on these
symbols with pairwise distinct unsigned indices.  Along the times
$n_\ell$ in Theorem~\ref{thm:walk}, the probability that any of the
requirements
\begin{align}
 |\tau(q(G_1,\ldots,G_m))|&<1/m
       &&(q\in\mathcal Q_m),\label{eq:selected-moments}\\
 \norm{D(v(G))-D(G_{i_1}^{\sigma_1})}_2&<1/m
       &&\bigl(v=x_{i_1}^{\sigma_1}\cdots x_{i_t}^{\sigma_t}
                   \in\mathcal V_m\bigr)\label{eq:selected-blocks}
\end{align}
fails tends to zero.  For the first line, this is a finite union of
the moment convergence statements in Theorem~\ref{thm:walk}.  For the
second line, the union bound and~\eqref{eq:independent-blocks} give
the upper bound $|\mathcal V_m|m^2e_{n_\ell}$.

Choose $\ell_m>\ell_{m-1}$, after fixing $m$, so large that
$n_{\ell_m}\ge N_m$ and this combined failure probability is less
than $c_m/2$.  Then the intersection of $E_{m,n_{\ell_m}}$ with all
the requirements~\eqref{eq:selected-moments}--\eqref{eq:selected-blocks}
has probability at least $c_m/2>0$.  Choose a deterministic tuple
\[
 (u_1^{(m)},\ldots,u_m^{(m)})\in\U(M)^m
\]
from that intersection.  The constants $c_m$ may tend to zero as fast
as they please: the choice of $\ell_m$ is made separately for each
fixed $m$, after the positive lower bound has been obtained.

Write $a_1=a$ and $a_{-1}=b$.  The endpoint conditions and
\eqref{eq:selected-blocks} imply, for every
$v=x_{i_1}^{\sigma_1}\cdots x_{i_t}^{\sigma_t}\in\mathcal V_m$,
\begin{equation}\label{eq:selected-multipliers}
 \norm{T(v(u^{(m)}))-a_{\sigma_1}v(u^{(m)})}_2
 =\norm{D(v(u^{(m)}))-a_{\sigma_1}}_2<2/m.
\end{equation}
The equality uses $T(v)=D(v)v$ and the $2$-norm isometry of right
multiplication by the unitary $v$.

\medskip\noindent
\emph{The exact family in the ultrapower.}
Fix a free ultrafilter $\omega_1$ on the tuple index $m$.  Set
$u_i^{(m)}=1$ when $m<i$, and define
\[
 N=M^{\omega_1},\qquad
 g_i=[(u_i^{(m)})_{m\ge1}]_{\omega_1}\quad(i\ge1),
 \qquad R=T^{\omega_1}.
\]
Lemma~\ref{lem:ultrapower} makes $R$ a bounded linear map on $N$,
with $\norm{R}\le C$, and supplies its bounded-ball $2$-continuity.
Each $g_i$ is unitary.  Every fixed nonidentity reduced word $q$ lies
in $\mathcal Q_m$ for all sufficiently large $m$, so
\eqref{eq:selected-moments} gives
\begin{equation}\label{eq:exact-haar-family}
 \tau^{\omega_1}(q(g_1,g_2,\ldots))=0.
\end{equation}
These are precisely the joint moments of a countable free Haar
family: powers of one generator are centered, and expansion of
alternating centered Laurent polynomials in different generators
reduces to nonidentity words.  The corresponding freeness of their
generated von Neumann algebras follows by bounded $2$-norm
approximation.

Identify $a,b\in M$ with their constant images in $N$.  For every
nonempty word $w=g_{i_1}^{\sigma_1}\cdots g_{i_t}^{\sigma_t}$ with
pairwise distinct unsigned indices,~\eqref{eq:selected-multipliers}
holds for all sufficiently large tuple indices.  Hence it becomes the
exact identity
\begin{equation}\label{eq:first-letter-multipliers}
 R(w)=
 \begin{cases}
  aw,&\sigma_1=1,\\
  bw,&\sigma_1=-1.
 \end{cases}
\end{equation}
These are all the hypotheses of Theorem~\ref{thm:rigidity}, applied
to $(N,\tau^{\omega_1})$ and $R$.  That theorem yields $a=b$ in $N$.
The constant embedding preserves the faithful trace, so this also
means $a=b$ in $M$, contrary to their choice.  Thus $Z_\infty$ is
deterministic, and the earlier argument proves
\eqref{eq:liouville-conclusion}.
\end{proof}

Only the scalar joint moments of the $g_i$ and the identities
\eqref{eq:first-letter-multipliers} enter the rigidity theorem.
Freeness of these unitaries from $a$ or $b$ is not asserted or needed.

\section{Rigidity of the first-letter identities}\label{sec:rigidity}

The identities furnished by the walk concern words whose unsigned indices
are pairwise distinct.  We prove that their two left multipliers must
agree.  This rules out the distinct martingale endpoints constructed
in Section~\ref{sec:liouville}.

\begin{theorem}\label{thm:rigidity}
Let $(N,\tau)$ be a finite von Neumann algebra with a faithful normal
tracial state.  Let $R:N\to N$ be bounded and complex linear, and suppose
that
\begin{equation}\label{eq:rigidity-continuity}
 \norm{R(x_n)}_2\longrightarrow0
 \quad\text{whenever}\quad
 \sup_n\norm{x_n}<\infty,
 \qquad \norm{x_n}_2\longrightarrow0.
\end{equation}
Let $(g_j)_{j\geq1}$ be a free Haar family in $N$, and let $a,b\in N$.
Suppose that for every nonempty signed word
$w=g_{i_1}^{\varepsilon_1}\cdots g_{i_q}^{\varepsilon_q}$ with
$\varepsilon_j\in\{1,-1\}$ and pairwise distinct indices $i_1,\ldots,i_q$,
\begin{equation}\label{eq:rigidity-first-letter}
 R(w)=
 \begin{cases}
 aw,&\varepsilon_1=1,\\
 bw,&\varepsilon_1=-1.
 \end{cases}
\end{equation}
Then $a=b$.
\end{theorem}

Here the free Haar condition means that the trace of every nonidentity
reduced group word in the $g_j$ is zero.  In particular, distinct reduced
words are orthonormal for the tracial $2$-norm.  The coefficients $a,b$
are arbitrary elements of the containing algebra.

The target is a unitary whose positive powers carry the multiplier $a$
and whose negative powers carry the multiplier $b$.  A spectral argument
will show that these identities force $a=b$ for a bounded map in a finite
tracial algebra.  Higher powers of any one $g_j$ repeat its unsigned index, so
the hypothesis does not give this target directly.  Instead, we average
many free generators, control the repeated indices in suitable polynomial
expressions, and pass to the polar part of the resulting element in a
further tracial ultrapower.

\subsection{Norm control for free words}

We use the following form of Haagerup's length estimate
\cite[Lemma~1.4]{Haagerup1979}.  We give the proof, including the transfer
from the regular representation to the containing tracial algebra.

\begin{lemma}\label{lem:rigidity-length}
Let $\mathbb F_r$ be the free group on $r$ generators and let $\lambda$
be its left regular representation.  For a finitely supported family
$(c_w)_{|w|=q}$ on the reduced words of length $q\geq0$,
\begin{equation}\label{eq:rigidity-length}
 \left\|\sum_{|w|=q}c_w\lambda(w)\right\|
 \leq(q+1)\left(\sum_{|w|=q}|c_w|^2\right)^{1/2}.
\end{equation}
If $g_1,\ldots,g_r$ have free Haar joint moments in a faithfully tracial
von Neumann algebra, every scalar-coefficient word polynomial in them
has the same norm as the corresponding polynomial in the regular
generators.  Consequently, \eqref{eq:rigidity-length} also holds with
$\lambda(w)$ replaced by $w(g)$, and its right-hand side is
$(q+1)\norm{\sum_{|w|=q}c_w w(g)}_2$.
\end{lemma}

\begin{proof}
Write $L=\sum_{|w|=q}c_w\lambda(w)$.  On the standard basis of
$\ell^2(\mathbb F_r)$, separate each product $wx$ according to its
exact number $t\in\{0,\ldots,q\}$ of cancellations at the joining
boundary.  This gives $L=\sum_{t=0}^q L_t$ on finitely supported vectors.
Fix $t$.  Such a product has the form
\[
 w=uv,\qquad x=v^{-1}h,\qquad wx=uh,
 \qquad |u|=q-t,\quad |v|=t,
\]
where all three displayed concatenations $uv$, $v^{-1}h$, and $uh$
are reduced.  Conversely, these conditions describe exactly the
products with $t$ cancellations, allowing any of $u,v,h$ to be empty
when its required length permits this.  In particular, $h$ may be empty
when the input is exhausted.

For this fixed $t$, an input basis word of length at least $t$ has a
unique prefix of length $t$ and hence a unique suffix $h$.  An output
basis word of length at least $q-t$ likewise determines $h$ uniquely by
removing its prefix of length $q-t$.  Thus the input spaces indexed by
different $h$ are orthogonal, as are the corresponding output spaces.
Relative to these two decompositions, $L_t$ is a direct sum of matrices
$B_h$, with columns indexed by $v$ and rows by $u$.  An entry is
$c_{uv}$ when the three reduction conditions hold and is zero otherwise.
The fixed cutting length makes $(u,v)\mapsto uv$ injective within each
matrix.  Therefore
\[
 \norm{B_h}\leq\norm{B_h}_{\mathrm{HS}}
 \leq\left(\sum_{|w|=q}|c_w|^2\right)^{1/2}.
\]
Taking the direct-sum norm bounds $L_t$ by the same quantity.  Summing
over the $q+1$ possible values of $t$ proves
\eqref{eq:rigidity-length}.

For the norm transfer, recall that if $A\geq0$ belongs to a von Neumann
algebra with faithful tracial state $\rho$, then
\begin{equation}\label{eq:rigidity-norm-moments}
 \norm{A}=\lim_{n\to\infty}\rho(A^n)^{1/n}.
\end{equation}
The upper bound is immediate.  If $0<c<\norm{A}$, the spectral
projection $e=1_{(c,\infty)}(A)$ is nonzero and has positive trace, so
$\rho(A^n)\geq c^n\rho(e)$.  Taking roots and then letting
$c\uparrow\norm{A}$ proves the formula.  For a word polynomial $p$,
the free Haar moments give, for every $n\geq1$,
\[
 \tau\bigl((p(g)^*p(g))^n\bigr)
 =\tau_{\mathbb F_r}\bigl((p(\lambda)^*p(\lambda))^n\bigr).
\]
Both traces are faithful; here $\tau_{\mathbb F_r}$ is the canonical
trace on the free group von Neumann algebra.  Applying
\eqref{eq:rigidity-norm-moments} gives $\norm{p(g)}=\norm{p(\lambda)}$.
The asserted $2$-norm formula follows from orthonormality of reduced
words.
\end{proof}

Under the hypotheses of Theorem~\ref{thm:rigidity}, set
\begin{equation}\label{eq:rigidity-normalized-sums}
 s_r=r^{-1/2}\sum_{j=1}^r g_j\in N.
\end{equation}
Lemma~\ref{lem:rigidity-length} gives $\norm{s_r}_2=1$ and
$\norm{s_r}\leq2$ for every $r$.  Fix a free ultrafilter $\omega_2$ on
$\N$ and put
\begin{equation}\label{eq:rigidity-second-ultrapower}
 \widehat N=N^{\omega_2},\qquad
 \widehat\tau=\tau^{\omega_2},\qquad
 s=(s_r)_{\omega_2}\in\widehat N,\qquad S=R^{\omega_2}.
\end{equation}
Lemma~\ref{lem:ultrapower} ensures that $S$ is bounded and retains
$2$-norm continuity on operator norm balls.  We identify $a,b$ with
their constant images in $\widehat N$.

In the application from Section~\ref{sec:liouville},
$N=M^{\omega_1}$ and $\widehat N=(M^{\omega_1})^{\omega_2}$.
The first index $m$ selects increasingly accurate unitary tuples; the
second index $r$ increases the number of exact free generators being
averaged in $N$, where the bound $\norm{s_r}\leq2$ is already available.
The hypotheses and notation of Theorem~\ref{thm:rigidity} remain in
force through the next two subsections.

Figure~\ref{fig:proof-map} locates the two ultrapowers in the Liouville
argument and records the remaining passage from these bounded sums to a
polar unitary and then to $a=b$.
\begin{figure}[htbp]
\centering
\small
\begin{tabular}{c}
\fbox{\parbox{0.86\linewidth}{\centering
The fixed walk has asymptotically free moments.\\
A hypothetical nonconstant harmonic-map martingale has\\
distinct support points $a,b\in M$.}}\\[3pt]
$\downarrow$\quad finite selection and the first ultrapower\\[3pt]
\fbox{\parbox{0.86\linewidth}{\centering
In $N=M^{\omega_1}$: exact free Haar unitaries $g_1,g_2,\ldots$\\
and exact first-letter identities for $R=T^{\omega_1}$.}}\\[3pt]
$\downarrow$\quad free-group length estimate in $N$\\[3pt]
\fbox{\parbox{0.86\linewidth}{\centering
$s_r=r^{-1/2}\sum_{j=1}^r g_j$ satisfies $\norm{s_r}\leq2$.\\
In each fixed product of blocks $s_r(s_r^*s_r)^d$, $d\geq0$,\\
terms with repeated surviving indices have total norm tending to zero.}}\\[3pt]
$\downarrow$\quad the second ultrapower and polar approximation\\[3pt]
\fbox{\parbox{0.86\linewidth}{\centering
In $\widehat N=N^{\omega_2}$: $s=(s_r)_{\omega_2}$ has a unitary
polar part $V$,\\
and $S=R^{\omega_2}$ satisfies $S(V^k)=aV^k$, $S(V^{-k})=bV^{-k}$.}}\\[3pt]
$\downarrow$\quad bounded sine sums and the finite faithful trace\\[3pt]
\fbox{\parbox{0.86\linewidth}{\centering
$a=b$, a contradiction. Hence $T(x)=T(1)x$.\\
Last-input averaging then produces the tracial cocycle primitive.}}
\end{tabular}
\caption{The Liouville contradiction and its use for cocycles.
The first ultrapower uses scalar moment convergence; the norm bound
is proved for the exact free family before the second ultrapower.
The constants $a,b$ are carried through the trace-preserving constant
embeddings, and the power identities hold for every $k\geq1$.}
\label{fig:proof-map}
\end{figure}


\subsection{From distinct indices to ordered blocks}

Polynomial expressions in $s_r$ may contain repeated indices, so their
multiplier identities do not follow directly from the hypothesis.
For polar approximation we will need powers of $sp(s^*s)$, where $p$
is a scalar polynomial.  Expanding these powers without commuting any
factors gives products of the blocks $B_d(z)=z(z^*z)^d$, $d\geq0$.
We establish their identities by showing that repeated surviving indices
contribute an error tending to zero in operator norm.

\begin{lemma}\label{lem:rigidity-blocks}
For every $k\geq1$ and $d_1,\ldots,d_k\geq0$, let
\[
 P(z)=B_{d_1}(z)B_{d_2}(z)\cdots B_{d_k}(z).
\]
Then
\begin{equation}\label{eq:rigidity-block-identities}
 S(P(s))=aP(s),\qquad S(P(s)^*)=bP(s)^*.
\end{equation}
\end{lemma}

\begin{proof}
Fix this product, and let $L=\sum_{j=1}^k(2d_j+1)$ be its degree.
Every term in the expansion of one block $B_d(s_r)$ has the odd
alternating sign pattern
\[
 +,-,+,-,\ldots,+.
\]
Deleting adjacent inverse letters deletes two consecutive opposite
signs.  The remaining pattern is still alternating, begins and ends
positively, and is nonempty: the excess of positive over negative
signs remains one.  Hence the free reduction of each block is a
nonempty word with positive endpoint signs.  After reducing the blocks
internally, their boundaries join two positive letters and admit no
further cancellation.  Every reduced word in $P(s_r)$ is therefore
nonempty and begins and ends positively.  Taking adjoints reverses
the block order and all signs, so the corresponding assertion for
$P(s_r)^*$ has negative endpoint signs.

We next bound the coefficient of every possible reduced word, including
words much shorter than $L$.  Write
\[
 P(s_r)=\sum_{q=1}^L\sum_{|w|=q}c_w(r)w(g).
\]
For a fixed reduced word $w$ of length $q$, record a free reduction of
an original term by pairing its original positions whenever their
letters cancel.  For definiteness, the reduction may be performed by
reading from left to right and canceling against the last surviving
letter whenever possible.  The recorded pairs form a noncrossing
partial matching of the $L$ positions: there are $(L-q)/2$ pairs, and
the unmatched positions spell $w$ in order.  All nested and newly
adjacent cancellations are included in this record.

There are at most $3^L$ noncrossing partial matchings of $L$ positions.
Indeed, mark each position as an opening endpoint, a closing endpoint,
or unmatched; noncrossing forces a closing endpoint to pair with the
last unpaired opening endpoint, so this encoding is injective.  We
may count every such partial matching, including ones not realized by
a reduction.  For a fixed pattern and fixed $w$, the labels at the
unmatched positions are prescribed.  Each canceled pair has equal
unsigned labels and opposite signs and hence has at most $r$ label
choices.  Coincidences between different pairs or with unmatched
positions can only reduce the number of choices.  Each original term
has coefficient $r^{-L/2}$.  Consequently, for every $q$ and $w$,
\begin{equation}\label{eq:rigidity-coefficient-bound}
 |c_w(r)|\leq C_L r^{(L-q)/2}r^{-L/2}
 =C_Lr^{-q/2},\qquad C_L=3^L.
\end{equation}
An impossible parity or reduction pattern contributes zero.  The same
bound applies to the adjoint polynomial.

Let $E_{r,q}$ be the sum of those length-$q$ terms whose \emph{surviving
reduced word} has a repeated unsigned index.  Original repetitions
which disappear in the reduction remain included in the coefficients
of the retained words.  For $q=1$ this sum is zero.  For $q\geq2$ the
number of possible words in the sum is at most
\[
 2^q\binom q2 r^{q-1}:
\]
choose a pair of equal-index positions, choose their common label and
the other $q-2$ labels, and choose all signs.  Counting nonreduced
strings here only enlarges the bound.  By orthonormality and
\eqref{eq:rigidity-coefficient-bound},
\begin{equation}\label{eq:rigidity-discard-ltwo}
 \norm{E_{r,q}}_2^2
 \leq C_L^2 2^q\binom q2 r^{-1}\qquad(2\leq q\leq L).
\end{equation}
Lemma~\ref{lem:rigidity-length} now yields the aggregate operator norm
estimate
\begin{equation}\label{eq:rigidity-discard-norm}
 E_r=\sum_{q=1}^L E_{r,q},\qquad
 \norm{E_r}\leq\sum_{q=2}^L(q+1)\norm{E_{r,q}}_2
 =O_L(r^{-1/2})\longrightarrow0.
\end{equation}
There is no scalar term, by the block reduction argument.

Put $Q_r=P(s_r)-E_r$.  Each word in this finite sum has pairwise
distinct unsigned indices and begins positively.  Thus
\eqref{eq:rigidity-first-letter} and linearity give $R(Q_r)=aQ_r$.
Moreover, $\norm{P(s_r)}\leq2^L$, so \eqref{eq:rigidity-discard-norm}
also proves that the sequence $(Q_r)$ is bounded.  In $N$ we have
\[
 \norm{R(P(s_r))-aP(s_r)}
 \leq(\norm{R}+\norm{a})\norm{E_r}\longrightarrow0.
\]
Passing to $\widehat N$ proves the first identity in
\eqref{eq:rigidity-block-identities}.  For the adjoint, the retained
sum is $Q_r^*$: its reduced words have distinct unsigned indices and
begin negatively.  Therefore $R(Q_r^*)=bQ_r^*$, and the identical
estimate with $b$ proves the second identity.
\end{proof}

\subsection{A polar unitary and all its powers}

The spectral calculation below is related to the polar decomposition
of circular elements in free probability, established by Voiculescu
and also proved by Banica~\cite{Voiculescu1990Circular,Banica1996Polar}.
We derive the required Catalan moments and spectral distribution
directly. Only the absence of an atom at zero is needed to obtain
the unitary polar part.

\begin{lemma}\label{lem:rigidity-polar}
The polar part $V$ of $s$ is a unitary in $\widehat N$, and
\begin{equation}\label{eq:rigidity-power-identities}
 S(V^k)=aV^k,\qquad S(V^{-k})=bV^{-k}
 \qquad(k\geq1).
\end{equation}
\end{lemma}

\begin{proof}
We first show that $s$ has zero kernel, which makes its polar part
unitary in the finite algebra.  We then transfer the block identities
to powers of that unitary.
Set $A=s^*s$, so $0\leq A\leq4$.  To determine its kernel, we compute
its moments.  For
$d\geq1$, the quantity $\tau((s_r^*s_r)^d)$ is $r^{-d}$ times the
number of assignments of labels in $\{1,\ldots,r\}$ to the alternating
$2d$ positions
\[
 -,+,-,+,\ldots,-,+
\]
whose word reduces to the identity.  An identity reduction pairs all
positions noncrossingly, so it uses at most $d$ unsigned labels.
Assignments using at most $d-1$ labels number $O_d(r^{d-1})$: for each
number of labels, choose an ordered list of those labels and then
assign a member of that list to each of the $2d$ positions.

If exactly $d$ labels occur in an identity word, every label occurs
twice.  Its two occurrences must be paired in any cancellation
matching, so the matching is uniquely determined by the labeling.
Conversely, in a noncrossing perfect matching, every pair has an even
number of positions inside it, since those positions must be matched
among themselves.  The endpoints therefore have opposite parities,
and hence opposite signs in our alternating string.  Assigning
distinct labels to its $d$ pairs gives an identity word: erase
innermost pairs successively.  It follows that the number of leading
assignments is exactly
\[
 C_d(r)_d,\qquad (r)_d=r(r-1)\cdots(r-d+1),
\]
where $C_d$ is the number of noncrossing perfect matchings on $2d$
positions.  Choosing the partner of the first position gives
$C_0=1$ and $C_d=\sum_{j=0}^{d-1}C_jC_{d-1-j}$.  The generating
function thus satisfies $C(z)=1+zC(z)^2$; its solution with constant
term one gives
\[
 C_d=\frac1{d+1}\binom{2d}{d}.
\]
Consequently,
\begin{equation}\label{eq:rigidity-catalan-moments}
 \widehat\tau(A^d)
 =\lim_{r\to\omega_2}\tau((s_r^*s_r)^d)
 =\lim_{r\to\infty}r^{-d}
       \bigl(C_d(r)_d+O_d(r^{d-1})\bigr)
 =C_d.
\end{equation}
The zeroth moment is also $C_0=1$.

Consider the measure
\begin{equation}\label{eq:rigidity-mp-law}
 d\nu(x)=\frac1{2\pi}\sqrt{\frac{4-x}{x}}\,
                1_{(0,4)}(x)\,dx.
\end{equation}
For every integer $d\geq0$, substitution $x=4t$ and the beta integral
give
\[
 \begin{split}
 \int_0^4 x^d\,d\nu(x)
 &=\frac{4^{d+1}}{2\pi}
     \int_0^1 t^{d-1/2}(1-t)^{1/2}\,dt\\
 &=\frac{4^{d+1}}{2\pi}
     \frac{\Gamma(d+\tfrac12)\Gamma(\tfrac32)}{\Gamma(d+2)}
 =\frac{(2d)!}{d!(d+1)!}=C_d.
 \end{split}
\]
In particular, $\nu$ is a probability measure.  Both $\nu$ and the
spectral distribution of $A$ are supported in $[0,4]$.  Equality of
all their moments implies equality on polynomials and then on all
continuous functions by uniform polynomial approximation.  They are
therefore the same measure.  Since \eqref{eq:rigidity-mp-law} has no
atom at zero,
\[
 \widehat\tau(1_{\{0\}}(A))=0,
 \qquad 1_{\{0\}}(A)=0
\]
by faithfulness.  In the polar decomposition $s=VA^{1/2}$ this gives
$V^*V=1$.  Traciality and faithfulness then give
\[
 \widehat\tau(1-VV^*)=1-\widehat\tau(V^*V)=0,
 \qquad VV^*=1.
\]
Thus $V$ is a unitary.

It remains to transfer the identities.  For a fixed regularization
parameter we will use polynomial approximation in operator norm;
we will then remove the regularization using bounded-ball $2$-continuity.
For $\varepsilon>0$ put
\[
 x_\varepsilon=s(A+\varepsilon1)^{-1/2}
 =V\left(\frac{A}{A+\varepsilon1}\right)^{1/2}.
\]
Fix $k\geq1$.  A polynomial $p(t)=\sum_{d=0}^D c_dt^d$ has the
literal ordered expansion
\begin{equation}\label{eq:rigidity-ordered-expansion}
 [sp(A)]^k
 =\sum_{d_1,\ldots,d_k=0}^D c_{d_1}\cdots c_{d_k}
       [sA^{d_1}]\cdots[sA^{d_k}].
\end{equation}
Every summand is a product covered by Lemma~\ref{lem:rigidity-blocks}.
That lemma and complex linearity yield
\[
 S([sp(A)]^k)=a[sp(A)]^k,
 \qquad S(([sp(A)]^k)^*)=b([sp(A)]^k)^*.
\]
For the second identity one takes the adjoint of the expansion
\eqref{eq:rigidity-ordered-expansion}, conjugating its scalar
coefficients and reversing the block order, and applies the separate
adjoint identity in Lemma~\ref{lem:rigidity-blocks}.
Choose polynomials converging uniformly on $[0,4]$ to
$(t+\varepsilon)^{-1/2}$.  Functional calculus and boundedness of $S$
pass the last two identities to the norm limit, giving
\begin{equation}\label{eq:rigidity-regularized-powers}
 S(x_\varepsilon^k)=ax_\varepsilon^k,
 \qquad S((x_\varepsilon^k)^*)=b(x_\varepsilon^k)^*.
\end{equation}

The $x_\varepsilon$ are contractions, and their distance to $V$
satisfies
\[
 \norm{x_\varepsilon-V}_2^2
 =\int_0^4
   \left(\sqrt{\frac{x}{x+\varepsilon}}-1\right)^2\,d\nu(x)
 \longrightarrow0
\]
by dominated convergence and $\nu(\{0\})=0$.  Multiplication satisfies
$\norm{xyz}_2\leq\norm{x}\norm{y}_2\norm{z}$ in a tracial algebra.
The noncommutative telescoping identity
\[
 x_\varepsilon^k-V^k
 =\sum_{j=0}^{k-1}x_\varepsilon^j
                  (x_\varepsilon-V)V^{k-1-j}
\]
therefore shows that
$\norm{x_\varepsilon^k-V^k}_2\leq k\norm{x_\varepsilon-V}_2\to0$;
the same holds for their adjoints.  All these powers are contractions.
Using the bounded-ball $2$-norm continuity of $S$ in
\eqref{eq:rigidity-regularized-powers}, along for example
$\varepsilon=1/n$, proves \eqref{eq:rigidity-power-identities}.
\end{proof}

\subsection{The spectral conclusion}

The final step uses boundedness, the identities for powers, and the
finite faithful trace. No continuity assumption on the map is needed
in Lemma~\ref{lem:rigidity-unitary}.

\begin{lemma}\label{lem:rigidity-unitary}
Let $(\mathcal A,\rho)$ be a finite von Neumann algebra with faithful
normal tracial state, let $V\in\mathcal A$ be unitary, and let
$S:\mathcal A\to\mathcal A$ be bounded and complex linear.  If
$a,b\in\mathcal A$ satisfy
\[
 S(V^k)=aV^k,\qquad S(V^{-k})=bV^{-k}\qquad(k\geq1),
\]
then $a=b$.
\end{lemma}

\begin{proof}
We use analytic polynomials whose imaginary parts remain uniformly
bounded on the unit circle while their real parts grow near $1$.
Rotating these polynomials
will bound $a-b$ on each short spectral arc; the trace will then combine
the compression estimates without a loss depending on the number of arcs.
Write $\mathbb T=\{z\in\C:|z|=1\}$ and
\[
 H_m(z)=\sum_{k=1}^m\frac{z^k}{k},\qquad
 h_m=\sum_{k=1}^m\frac1k.
\]
We first verify the uniform sine-sum bound
\begin{equation}\label{eq:rigidity-sine-bound}
 \sup_{m\geq1}\sup_{\theta\in\R}
 \left|\sum_{k=1}^m\frac{\sin(k\theta)}k\right|
 \leq B,\qquad B=1+\pi.
\end{equation}
By periodicity and oddness it suffices to use the angular distance
$t\in[0,\pi]$ from $2\pi\Z$.  At $t=0$ the sum is zero.  For $t>0$
put $K=\lfloor1/t\rfloor$.  The part with $k\leq\min(K,m)$ has
absolute sum at most $Kt\leq1$, since $|\sin(kt)|\leq kt$.
Geometric summation gives, for every interval of integers $p\leq q$,
\[
 \left|\sum_{k=p}^q\sin(kt)\right|
 \leq\left|\sum_{k=p}^q e^{ikt}\right|
 \leq\frac1{\sin(t/2)}\leq\frac\pi t.
\]
If $m\geq K+1$, put $p=K+1$ and
$A_j=\sum_{k=p}^j\sin(kt)$.  Abel summation gives
\[
 \sum_{k=p}^m\frac{\sin(kt)}k
 =\frac{A_m}{m}
   +\sum_{j=p}^{m-1}A_j\left(\frac1j-\frac1{j+1}\right).
\]
Its absolute value is at most $\pi/(tp)\leq\pi$, including when
$K=0$.  This proves \eqref{eq:rigidity-sine-bound}.

For $\lambda\in\mathbb T$, write
\[
 H_m(\lambda V)=X_{m,\lambda}+iY_{m,\lambda},
 \qquad X_{m,\lambda}=X_{m,\lambda}^*,\quad
 Y_{m,\lambda}=Y_{m,\lambda}^*.
\]
Continuous functional calculus and \eqref{eq:rigidity-sine-bound}
give $\norm{Y_{m,\lambda}}\leq B$.  The identities for positive and
negative powers imply
\[
 \begin{split}
 S\bigl(H_m(\lambda V)-H_m(\lambda V)^*\bigr)
 &=aH_m(\lambda V)-bH_m(\lambda V)^*\\
 &=(a-b)X_{m,\lambda}+i(a+b)Y_{m,\lambda}.
 \end{split}
\]
Since the argument of $S$ has norm at most $2B$, we obtain
\begin{equation}\label{eq:rigidity-real-bound}
 \norm{(a-b)X_{m,\lambda}}\leq C,
 \qquad C=B(2\norm{S}+\norm{a}+\norm{b}),
\end{equation}
uniformly in $m$ and $\lambda$.

Fix $m$ and any half-open arc $I\subset\mathbb T$ of angular length
at most $1/m$, and let $e=1_I(V)$.  Choose $\lambda$ to rotate the
midpoint of $I$ to $1$.  On the rotated arc the argument $\theta$
satisfies $|\theta|\leq1/m$, so
\[
 \operatorname{Re}H_m(e^{i\theta})
 =\sum_{k=1}^m\frac{\cos(k\theta)}k
 \geq\cos(1)h_m.
\]
Thus, for $X=X_{m,\lambda}$, Borel functional calculus gives an
element $F\in W^*(V)$ equal to $1/X$ on $I$ and to zero off $I$,
with
\[
 XF=e,\qquad
 \norm{F}\leq\frac1{\cos(1)h_m}.
\]
Writing $c=a-b$ and multiplying \eqref{eq:rigidity-real-bound} on
the right by $F$ gives
\begin{equation}\label{eq:rigidity-arc-bound}
 \norm{ce}=\norm{cXF}
 \leq\frac{C}{\cos(1)h_m}=:\delta_m.
\end{equation}
In the corner $e\mathcal A e$, this says
$ec^*ce\leq\delta_m^2 e$.

Partition $\mathbb T$ into finitely many disjoint half-open arcs of
angular length at most $1/m$, and denote their spectral projections
by $e_1,\ldots,e_{J_m}$.  They are orthogonal and sum to $1$.
Traciality and \eqref{eq:rigidity-arc-bound} imply
\[
 \begin{split}
 \norm{c}_2^2
 &=\sum_{j=1}^{J_m}\rho(c^*ce_j)
  =\sum_{j=1}^{J_m}\rho(e_jc^*ce_j)\\
 &\leq\delta_m^2\sum_{j=1}^{J_m}\rho(e_j)
  =\delta_m^2\longrightarrow0,
 \end{split}
\]
because $h_m\geq\log(m+1)\to\infty$.  Each arc estimate is weighted
by the trace of its own projection; the number of arcs introduces no
factor.  Faithfulness gives $c=0$, as required.  The right multiplication
by $F$ and the tracial compression identity apply to arbitrary $a,b$;
the argument imposes no commutation condition and uses no particular
spectral distribution of $V$.
\end{proof}

\begin{proof}[Proof of Theorem~\ref{thm:rigidity}]
Construct $s\in\widehat N$ and $S=R^{\omega_2}$ as in
\eqref{eq:rigidity-second-ultrapower}.
Lemma~\ref{lem:rigidity-blocks} propagates the first-letter identities
to every ordered block product.  Lemma~\ref{lem:rigidity-polar} then
provides a unitary $V$ satisfying the positive and negative power
identities.  Lemma~\ref{lem:rigidity-unitary} gives equality of the
constant images of $a$ and $b$ in $\widehat N$.  The constant embedding
is trace preserving and faithful, so $a=b$ in $N$.
\end{proof}

\begin{remark}\label{rem:finite-trace-essential}
The finite hypothesis in Lemma~\ref{lem:rigidity-unitary} is essential.
On $H=\ell^2(\Z)$ let $Ve_n=e_{n+1}$, let $p$ project onto $\C e_0$,
and let $q$ project onto the closed span of the $e_n$ with $n<0$.
The contraction $S(X)=pXq$ on $B(H)$ satisfies
$S(V^k)=pV^k$ and $S(V^{-k})=0$ for every $k\ge1$:
the only nonzero input coordinate of $pV^k$ is $e_{-k}$, whereas
that of $pV^{-k}$ is $e_k$.
Thus the power identities hold with $a=p\ne0=b$.
\end{remark}

\section{Bounded primitives for Hochschild cocycles}\label{sec:cohomology}

We first apply Theorem~\ref{thm:liouville} to separately normal
cocycles on a type $\mathrm{II}_1$ algebra with separable predual and a
faithful normal tracial state.  The resulting uniform bound for a
primitive will allow us first to remove separability and then to
combine central summands.  Throughout this
section, normality of a linear map means ultraweak continuity, and
separate normality has the corresponding meaning for multilinear maps.

\subsection{Normality and continuity on bounded balls}

The continuity needed by Theorem~\ref{thm:liouville} is the tracial
form of Akemann's continuity theorem
\cite[Corollary~III.10]{Akemann1967}. We include a direct proof that
allows a center and does not require a separable predual.
The stronger joint multilinear continuity result for type
$\mathrm{II}_1$ factors is \cite[Theorem~4.4]{CPSS2003}.

\begin{lemma}\label{lem:normal-two-continuity}
Let $M$ be a finite von Neumann algebra with faithful normal tracial
state $\tau$, and let $B:M\to M$ be bounded and normal.  For every
$R<\infty$ and $\varepsilon>0$ there is $\delta>0$ such that
\[
 \norm{x}\le R,\quad \norm{x}_2<\delta
 \quad\Longrightarrow\quad \norm{B(x)}_2<\varepsilon.
\]
Consequently $B$ is $2$-norm continuous on every operator norm ball.
\end{lemma}

\begin{proof}
We argue by contradiction.  A failure of continuity will give elements
with mutually orthogonal two-sided supports whose images stay uniformly
away from zero in $2$-norm.  Random signs will then contradict the
boundedness of $B$.

If the first assertion fails, there are $R>0$, $\eta>0$, and a sequence
$(x_j)$ such that
\[
 \norm{x_j}\le R,\qquad \norm{x_j}_2\longrightarrow0,
 \qquad \norm{B(x_j)}_2\ge\eta.
\]
Choose $\alpha>0$ with $\norm{B}\alpha<\eta/4$.  In the polar
decomposition $x_j=v_j|x_j|$, set
\[
 r_j=1_{(\alpha,\infty)}(|x_j|),\qquad
 y_j=v_j|x_j|r_j,\qquad
 \ell_j=v_jr_jv_j^*,\qquad e_j=r_j\vee\ell_j.
\]
Here $r_j$ and $\ell_j$ are the right and left support projections of
$y_j$.  In particular,
\[
 y_j=e_jy_je_j,\quad \norm{y_j}\le R,\quad
 \norm{x_j-y_j}\le\alpha,\quad
 \norm{B(y_j)}_2\ge 3\eta/4.
\]
The trace and the spectral truncation give
\[
 \tau(e_j)\le\tau(r_j)+\tau(\ell_j)
 =2\tau(r_j)\le2\alpha^{-2}\norm{x_j}_2^2.
\]
Pass to a subsequence for which $\sum_j\tau(e_j)<\infty$, and put
$p_j=\bigvee_{h\ge j}e_h$.  Then $p_j$ decreases to zero: indeed,
$\tau(p_j)\le\sum_{h\ge j}\tau(e_h)\to0$, and $\tau$ is faithful
and normal.

For fixed $j$ and $h>j$, the compressions
\[
 (p_j-p_h)y_j(p_j-p_h)
\]
converge ultraweakly to $y_j$ as $h\to\infty$.  Applying $B$ preserves
this convergence.  The $2$-norm is ultraweakly lower semicontinuous,
as is seen from
\[
 \norm{y}_2=
 \sup\{ |\tau(b^*y)|:b\in M,\ \norm{b}_2\le1\}.
\]
Starting with $j_1=1$, we may therefore choose successively
$j_{s+1}>j_s$ so that
\[
 t_s=(p_{j_s}-p_{j_{s+1}})y_{j_s}(p_{j_s}-p_{j_{s+1}}),
 \qquad \norm{B(t_s)}_2\ge\eta/2.
\]
The $t_s$ have pairwise orthogonal two-sided supports and
$\norm{t_s}\le R$.  Hence
$\norm{\sum_{s=1}^N\epsilon_st_s}\le R$ for every choice of signs
$\epsilon_s\in\{-1,1\}$.  Taking independent uniform signs, Hilbert
space orthogonality of these scalar random variables yields
\[
 \frac{N\eta^2}{4}
 \le\sum_{s=1}^N\norm{B(t_s)}_2^2
 =\E\norm{B\Bigl(\sum_{s=1}^N\epsilon_st_s\Bigr)}_2^2
 \le\norm{B}^2R^2,
\]
which is impossible for large $N$.  This proves the asserted modulus.
Applying it to differences, which lie in the ball of radius $2R$,
gives continuity on the ball of radius $R$.
\end{proof}

\subsection{Averaging the last input}

\begin{proposition}\label{prop:separable-primitive}
Let $M$ be of type $\mathrm{II}_1$ with separable predual and faithful
normal tracial state $\tau$.  If $k\ge2$ and
$f\in C_b^k(M,M)$ is a separately normal cocycle, then there is
$g\in C_b^{k-1}(M,M)$ such that
\begin{equation}\label{eq:separable-primitive-bound}
 dg=f,\qquad \norm{g}\le\norm{f}.
\end{equation}
\end{proposition}

\begin{proof}
We first construct an averaged cochain that is a right module map in
its last input.  Evaluation at $1$ will give a $(k-1)$-cochain, and averaging
the original equation $df=0$ will identify its differential.

Choose the symmetric countably supported measure $\mu$ for which
Theorem~\ref{thm:liouville} applies, and write
$\mu=\sum_{s\ge1}p_s\delta_{u_s}$.  For a bounded $k$-cochain
$\Phi$, define
\begin{equation}\label{eq:last-input-average}
 (P\Phi)(\mathbf a,x)
 =\sum_{s\ge1}p_s\Phi(\mathbf a,xu_s)u_s^*,
 \qquad \mathbf a=(a_1,\ldots,a_{k-1}).
\end{equation}
The sum converges in operator norm, and $P$ is a contraction for the
cochain norm.  If $W_j=U_1\cdots U_j$, where the $U_i$ are independent
with law $\mu$, and $W_0=1$, induction gives
\[
 (P^j f)(\mathbf a,x)
 =\E\bigl[f(\mathbf a,xW_j)W_j^*\bigr].
\]
Set
\[
 F_n=\frac1n\sum_{j=0}^{n-1}P^j f.
\]

We use compactness for pointwise ultraweak convergence.  For each fixed
tuple $\mathbf x=(x_1,\ldots,x_k)$, all values $F_n(\mathbf x)$ lie
in the ultraweakly compact ball of radius
$\norm{f}\prod_i\norm{x_i}$.  The product of these balls, indexed by
all tuples in $M^k$, is compact.  Multilinearity is closed for this
topology, since addition and scalar multiplication are ultraweakly
continuous.  Thus some subnet of $(F_n)$ converges at every tuple to
a bounded multilinear map $F$ with $\norm{F}\le\norm{f}$.

On any common cochain norm ball, $P$ is continuous for this topology.
To see this, suppose $\Phi_\lambda\to\Phi$ pointwise ultraweakly and
$\norm{\Phi_\lambda},\norm{\Phi}\le C$.  For a fixed tuple
$(\mathbf a,x)$ and a normal functional $\rho\in M_*$, truncate
\eqref{eq:last-input-average} after $L$ terms.  The finite sum passes
to the limit, and the remaining difference has absolute value at most
\[
 2C\norm{\rho}\norm{x}
 \prod_{i=1}^{k-1}\norm{a_i}\sum_{s>L}p_s.
\]
Letting $L\to\infty$ proves the claim.  Since
\[
 PF_n-F_n=\frac{P^nf-f}{n},\qquad
 \norm{PF_n-F_n}\le\frac{2\norm{f}}n,
\]
the subnet limit satisfies $PF=F$.

Fix $\mathbf a$, and put $B(x)=f(\mathbf a,x)$.  This is a bounded
normal linear map.  For any $R<\infty$ and $\varepsilon>0$,
Lemma~\ref{lem:normal-two-continuity} provides $\delta>0$ such that
\[
 \norm{x}\le R,\quad \norm{x}_2<\delta
 \quad\Longrightarrow\quad
 \norm{B(xv)v^*}_2<\varepsilon
 \quad\hbox{for every }v\in\U(M).
\]
Indeed, right multiplication by a unitary preserves both the operator
norm and the $2$-norm.  The same bound holds after averaging over any
$W_j$ and after taking the Ces\`aro average.  Ultraweak lower
semicontinuity of the $2$-norm then gives
$\norm{F(\mathbf a,x)}_2\le\varepsilon$ under these input bounds.
Using $\varepsilon/2$ in place of $\varepsilon$ proves the required
continuity modulus for $T(x)=F(\mathbf a,x)$.  The equality $PF=F$
is exactly the harmonicity identity for this map.  Consequently
Theorem~\ref{thm:liouville} gives
\begin{equation}\label{eq:averaged-right-module}
 F(\mathbf a,x)=h(\mathbf a)x,
 \qquad h(\mathbf a)=F(\mathbf a,1).
\end{equation}
In particular $h$ is a bounded $(k-1)$-cochain and
$\norm{h}\le\norm{f}$.

It remains to obtain its coboundary identity.  Let
$\nu_n=n^{-1}\sum_{j=0}^{n-1}\mu^{*j}$ and
$h_n(\mathbf a)=F_n(\mathbf a,1)$.  Evaluate the original equation
$df=0$ at $(a_1,\ldots,a_k,v)$, multiply the output on the right by
$v^*$, and average with respect to $v\sim\nu_n$.  This gives the
exact equality
\begin{align}
 0={}&a_1h_n(a_2,\ldots,a_k)
 +\sum_{j=1}^{k-1}(-1)^j
 h_n(a_1,\ldots,a_ja_{j+1},\ldots,a_k)\notag\\
 &+(-1)^k F_n(a_1,\ldots,a_{k-1},a_k)
 +(-1)^{k+1}f(a_1,\ldots,a_k).
 \label{eq:original-cocycle-average}
\end{align}
The final term is unchanged because $vv^*=1$.  Along the chosen
subnet, each displayed evaluation converges ultraweakly; multiplication
by a fixed $a_i$ preserves that convergence.  Passing to the limit
and using \eqref{eq:averaged-right-module} therefore yields
\[
 0=(dh)(a_1,\ldots,a_k)+(-1)^{k+1}f(a_1,\ldots,a_k).
\]
Thus $g=(-1)^kh$ satisfies \eqref{eq:separable-primitive-bound}.
This calculation averages $df$ at the displayed tuples; it requires
neither a commutation relation between $P$ and $d$ nor separate
normality of the limiting cochain $F$.
\end{proof}

\subsection{Removing the separability assumption}

Uniform primitive bounds and pointwise ultraweak compactness were
used to remove separability in \cite[Section~7]{CPSS2003}.
As in that argument, conditional expectations allow the approximating subalgebras
to be chosen independently for each finite set of inputs.

\begin{proposition}\label{prop:tracial-primitive}
Let $M$ be any type $\mathrm{II}_1$ von Neumann algebra admitting a
faithful normal tracial state.  Every separately normal cocycle
$f\in C_b^k(M,M)$, $k\ge2$, has a bounded primitive $g$ with
$\norm{g}\le\norm{f}$.
\end{proposition}

\begin{proof}
Order the finite subsets $E\subset M$ by inclusion.  For each $E$,
choose a countably generated unital von Neumann subalgebra $N_E$
containing $E$, the finite set $f(E^k)$, and unital matrix systems
of every size $m\ge2$.  Such systems exist in a type $\mathrm{II}_1$
algebra: the unit can be divided into $m$ equivalent orthogonal
projections~\cite[Examples~III.2.5.4(iii)]{Blackadar2006}, and partial isometries between them supply the matrix
units.  Adjoining one system for each $m$ adds only countably many
generators.

The restriction of the trace to $N_E$ is faithful and normal.  The
tracial Hilbert space $L^2(N_E)$ is separable: $*$-polynomials with
coefficients in $\mathbb Q+i\mathbb Q$ in a countable set of generators are dense in the
generated $C^*$-algebra in norm, and bounded strong approximation
makes that algebra $2$-norm dense in $N_E$.  The faithful tracial
representation thus gives $N_E$ separable predual.  Moreover $N_E$
has no type I central part.  A nonzero homogeneous type $\mathrm{I}_n$
central summand would inherit a unital copy of $M_m(\C)$ for every
$m$, whereas its $n$-dimensional fibers cannot admit a unital copy
when $m>n$.  Every finite type I algebra decomposes into its
homogeneous type $\mathrm{I}_n$ summands.  Since $N_E$ is finite,
it is therefore of type $\mathrm{II}_1$.

Let $\mathcal E_E:M\to N_E$ be the normal trace-preserving
conditional expectation. Its existence and bimodularity are classical;
see \cite[Th\'eor\`eme~8]{Dixmier1953} and
\cite[Theorem~1]{Umegaki1954}. Define
\[
 f_E=\mathcal E_E\circ f\big|_{N_E^k}.
\]
This map is separately normal and has norm at most $\norm{f}$.
For inputs in $N_E$, bimodularity of $\mathcal E_E$ gives
$df_E=\mathcal E_E(df)=0$.  By
Proposition~\ref{prop:separable-primitive}, choose
$g_E\in C_b^{k-1}(N_E,N_E)$ such that
$dg_E=f_E$ and $\norm{g_E}\le\norm{f}$.  Extend it to a cochain on
$M$ by
\[
 G_E(x_1,\ldots,x_{k-1})
 =g_E(\mathcal E_Ex_1,\ldots,\mathcal E_Ex_{k-1}).
\]
This extension has the same norm bound.  If
$(a_1,\ldots,a_k)\in E^k$, all its inputs and adjacent products
belong to $N_E$.  Expanding the differential consequently gives
\begin{equation}\label{eq:finite-set-exactness}
 (dG_E)(a_1,\ldots,a_k)
 =(dg_E)(a_1,\ldots,a_k)
 =\mathcal E_Ef(a_1,\ldots,a_k)
 =f(a_1,\ldots,a_k),
\end{equation}
where the final equality uses $f(E^k)\subset N_E$.

The common norm bound and the same product compactness used above
give a pointwise ultraweakly convergent subnet of $(G_E)$, with limit
$g\in C_b^{k-1}(M,M)$ and $\norm{g}\le\norm{f}$.  For a fixed
$k$-tuple $(a_1,\ldots,a_k)$, the indices of this subnet eventually
contain $E_0=\{a_1,\ldots,a_k\}$.  Equation~\eqref{eq:finite-set-exactness}
therefore holds eventually along the subnet.  Passing to the limit
in its finitely many coboundary terms gives $dg=f$ at this tuple.
This holds at every tuple, proving the proposition.  The choices of
$N_E$ and $\mathcal E_E$ can be made independently for different
$E$; eventual exactness is supplied by the directed finite subsets.
\end{proof}

\subsection{Central cuts and mixed inputs}

Uniform bounds and central decomposition also enter the tensor-product
argument of \cite[Section~4]{PopSmith2010}. In the present setting we
use an explicit homotopy to correct all inputs outside a given central
summand before assembling the primitives.

For a central projection $z\in M$, regard $zM$ as an $M$-bimodule
under multiplication and put $q=1-z$.  In degree zero we use
$C_b^0(M,zM)=zM$ and $(d\xi)(a)=a\xi-\xi a$.
On cochains with values in
$zM$, define the central cut by
\[
 (\operatorname{cut}_z\psi)(a_1,\ldots,a_n)
 =\psi(za_1,\ldots,za_n).
\]
In degree zero this means $\operatorname{cut}_z=\id$ on $zM$.
Centrality and the fact that $q$ annihilates this module show that
$\operatorname{cut}_z$ commutes with $d$.
For a cocycle $\psi$ with values in $zM$, a primitive for its
restriction to $zM$, pulled back to $M$, accounts only for
$\operatorname{cut}_z\psi$.
The following homotopy supplies the remaining difference
$\psi-\operatorname{cut}_z\psi$ on arbitrary inputs, with a bound
independent of $z$.

\begin{lemma}\label{lem:central-homotopy}
For $n\ge2$ and $\psi\in C_b^n(M,zM)$, set
\begin{equation}\label{eq:central-homotopy-formula}
 \begin{split}
 (J_z\psi)(a_1,\ldots,a_{n-1})
 =\sum_{i=1}^{n-1}(-1)^i\psi(&za_1,\ldots,za_{i-1},
 q,qa_i,\\[-2pt]
 &a_{i+1},\ldots,a_{n-1}).
 \end{split}
\end{equation}
Set $J_z=0$ in degrees zero and one.  Then, in every degree,
\begin{equation}\label{eq:central-homotopy-identity}
 dJ_z+J_zd=\id-\operatorname{cut}_z.
\end{equation}
For $n\ge1$ the degree-$n$ operator satisfies
\begin{equation}\label{eq:central-homotopy-bound}
 \norm{J_z\psi}\le(n-1)\norm{\psi}.
\end{equation}
\end{lemma}

\begin{proof}
The norm estimate follows directly from
\eqref{eq:central-homotopy-formula}.  To verify the identity, expand
each input as $a_i=za_i+qa_i$.  By multilinearity it suffices to
use tuples whose entries each belong to $zM$ or to $qM$.  On such
a tuple with first $qM$ entry at position $r$, formula
\eqref{eq:central-homotopy-formula} consists of the single term
inserting $q$ immediately before that entry, with sign $(-1)^r$;
on an all-$zM$ tuple it is zero.

Consider such an $n$-tuple with first $qM$ entry $a_r$.
In $(J_zd)\psi$, merging the inserted $q$ with $a_r$ has total
sign $(-1)^r(-1)^r=1$ and contributes
$\psi(a_1,\ldots,a_n)$.  A merger strictly before the insertion
has the same arguments as its counterpart in $(dJ_z)\psi$,
but the latter inserts at position $r-1$, so their signs are
opposite.  A merger strictly after the pair $(q,a_r)$ also has
the same arguments as its counterpart in $(dJ_z)\psi$, and its
merger index is larger by one, again reversing the sign.  The
merger of $a_{r-1}$ with the inserted $q$, when $r>1$, vanishes
because $zMqM=0$; the corresponding merger across $a_{r-1},a_r$
in $(dJ_z)\psi$ vanishes for the same reason.

The left outer terms cancel when $r>1$, with insertion indices
$r$ and $r-1$.  When $r=1$, both left outer terms vanish because
$qM$ annihilates the output module.  For the right outer terms,
if $r<n$ the same insertion index and the two consecutive
differential degrees give opposite signs.  If $r=n$, the last
input is in $qM$, so both right outer terms vanish.  Thus the
only surviving term is $\psi(a_1,\ldots,a_n)$, as required
because the cut of this tuple is zero.  On an all-$zM$ tuple,
both sides of \eqref{eq:central-homotopy-identity} vanish.

For completeness, in degree one the formula reads
\[
 (J_zd\psi)(a)=-(d\psi)(q,qa)=\psi(qa)
 =\psi(a)-\psi(za),
\]
since both outer module actions in $(d\psi)(q,qa)$ vanish.
In degree zero both sides are zero by the definitions.  This
proves the identity in all degrees.
\end{proof}

In particular, if $\psi\in C_b^k(M,zM)$ is a cocycle, then
\begin{equation}\label{eq:central-cocycle-correction}
 d(J_z\psi)=\psi-\operatorname{cut}_z\psi,
 \qquad \norm{J_z\psi}\le(k-1)\norm{\psi}.
\end{equation}
This bound is independent of the algebra and the central projection.

\subsection{Completion of the proof}

\begin{proof}[Proof of Theorem~\ref{thm:main}]
The zero algebra is immediate, so assume $M\ne0$.
Fix $k\ge2$ and an arbitrary bounded cocycle
$f_{\mathrm{orig}}\in C_b^k(M,M)$.  By the normal reduction in
Proposition~\ref{prop:classical}, there are a separately normal
cocycle $f$ and a bounded $(k-1)$-cochain $b$ such that
\[
 f_{\mathrm{orig}}=f+db.
\]
Let $z_{\mathrm{II}_1}$ be the central projection of the type
$\mathrm{II}_1$ summand, and let $z_0=1-z_{\mathrm{II}_1}$.
We first partition $z_{\mathrm{II}_1}$ into central projections
whose algebras admit faithful normal tracial states.

In the abelian algebra $Z(z_{\mathrm{II}_1}M)$, choose a maximal
orthogonal family $(z_i)_{i\in I}$ of nonzero supports of normal
states.  Its join is $z_{\mathrm{II}_1}$: a nonzero complementary
projection would support a further normal state, contradicting
maximality.  A state is faithful on its support, so each
$z_iZ(M)$ has a faithful normal state $\varphi_i$.  Composing
$\varphi_i$ with the faithful normal center-valued trace of $z_iM$
gives a faithful normal tracial state $\tau_i$ on $z_iM$;
see \cite[Theorem~III.2.5.7 and Corollary~III.2.5.8]{Blackadar2006}.
The family $I$ may be uncountable.  Omit $z_0$ if it is zero,
and otherwise keep the entire complementary summand as one piece.

For each piece define the restricted cocycle
\[
 f_i(x_1,\ldots,x_k)=z_i f(x_1,\ldots,x_k),
 \qquad x_j\in z_iM.
\]
Multiplication by a central projection is a bimodule map, so
$f_i$ is a separately normal cocycle and
$\norm{f_i}\le\norm{f}$.  For $i\in I$,
Proposition~\ref{prop:tracial-primitive} supplies a primitive
$b_i\in C_b^{k-1}(z_iM,z_iM)$ with
\[
 db_i=f_i,\qquad \norm{b_i}\le\norm{f}.
\]
If $z_0\ne0$, the vanishing result for algebras without a type
$\mathrm{II}_1$ central summand in
Proposition~\ref{prop:classical} supplies a bounded primitive
$b_0$ of $f_0$.  Its norm is some finite number.

Pull each primitive back to $M$ through the homomorphism
$x\mapsto z_ix$:
\[
 \widetilde b_i(x_1,\ldots,x_{k-1})
 =b_i(z_ix_1,\ldots,z_ix_{k-1}).
\]
The outer module actions also pull back, since the outputs lie
in $z_iM$.  Thus
\[
 d\widetilde b_i(x_1,\ldots,x_k)
 =z_i f(z_ix_1,\ldots,z_ix_k)
 =\operatorname{cut}_{z_i}(z_if)(x_1,\ldots,x_k).
\]
Define cochains with values in $z_iM$ on all of $M$ by
\[
 H_i=\widetilde b_i+J_{z_i}(z_if).
\]
Equation~\eqref{eq:central-cocycle-correction} gives
\[
 dH_i=z_if,
 \qquad
 \norm{H_i}\le k\norm{f}\quad(i\in I),
\]
and, when present,
\[
 \norm{H_0}\le\norm{b_0}+(k-1)\norm{f}.
\]
There is consequently a single finite constant $C$ bounding the
norms of all these cochains; one can take the maximum of the two
displayed bounds, omitting the second when $z_0=0$.

For an arbitrary central partition of unity, the map
\[
 M\longrightarrow\prod_i^{\ell^\infty}z_iM,
 \qquad x\longmapsto(z_ix)_i,
\]
is an isometric $*$-isomorphism onto the bounded direct product.
For surjectivity, the finite partial sums of a bounded family
$(x_i)_i$, $x_i\in z_iM$, form a net directed by inclusion of the
finite index sets.  This net converges strongly, since the central
supports are orthogonal; its limit belongs to $M$.
For each tuple $(x_1,\ldots,x_{k-1})$, the family
$(H_i(x_1,\ldots,x_{k-1}))_i$ belongs to this product and has norm
at most $C\prod_j\norm{x_j}$.  Let $G(x_1,\ldots,x_{k-1})\in M$
be the corresponding element.  Coordinatewise multilinearity
shows that $G$ is a bounded multilinear cochain with
$\norm{G}\le C$.  Central projection commutes with every term
of its coboundary, so
\[
 z_i(dG)=dH_i=z_if
 \qquad\hbox{for every central piece }i.
\]
Since the projections have join $1$, it follows that $dG=f$.
Finally,
\[
 d(G+b)=f+db=f_{\mathrm{orig}}.
\]
We have obtained an actual bounded $(k-1)$-cochain whose
coboundary is the given cocycle.  Hence
$H_b^k(M,M)=0$ for every complex von Neumann algebra $M$ and
every $k\ge2$, with the actual image of $d$ in the definition of
cohomology and with no complete-boundedness hypothesis.
\end{proof}


\clearpage
\begin{thebibliography}{CPSS03}
\raggedright

\bibitem[Ake67]{Akemann1967}
C.~A.~Akemann,
\emph{The dual space of an operator algebra},
Trans. Amer. Math. Soc. \textbf{126} (1967), no.~2, 286--302.
\href{https://doi.org/10.1090/S0002-9947-1967-0206732-8}{doi:10.1090/S0002-9947-1967-0206732-8}.

\bibitem[Ban96]{Banica1996Polar}
T.~Banica,
\emph{On the polar decomposition of circular variables},
Integral Equations Operator Theory \textbf{24} (1996), 372--377.
\href{https://doi.org/10.1007/BF01204607}{doi:10.1007/BF01204607}.

\bibitem[Bla06]{Blackadar2006}
B.~Blackadar,
\emph{Operator Algebras: Theory of $C^*$-Algebras and von Neumann Algebras},
Encyclopaedia of Mathematical Sciences, vol.~122,
Springer, Berlin--Heidelberg, 2006.
\href{https://bruceblackadar.com/Mathematics/Cycr.pdf}{Author's revised
online text dated February 8, 2017}.
\href{https://doi.org/10.1007/3-540-28517-2}{doi:10.1007/3-540-28517-2}.

\bibitem[Cam09]{Cameron2009}
J.~M.~Cameron,
\emph{Hochschild cohomology of $\mathrm{II}_1$ factors with Cartan
maximal abelian subalgebras},
Proc. Edinburgh Math. Soc. \textbf{52} (2009), no.~2, 287--295.
\href{https://doi.org/10.1017/S0013091507000053}{doi:10.1017/S0013091507000053}.

\bibitem[CES87]{CES1987}
E.~Christensen, E.~G.~Effros, and A.~M.~Sinclair,
\emph{Completely bounded multilinear maps and $C^*$-algebraic cohomology},
Invent. Math. \textbf{90} (1987), no.~2, 279--296.
\href{https://doi.org/10.1007/BF01388706}{doi:10.1007/BF01388706}.

\bibitem[CPSS03]{CPSS2003}
E.~Christensen, F.~Pop, A.~M.~Sinclair, and R.~R.~Smith,
\emph{Hochschild cohomology of factors with property $\Gamma$},
Ann. of Math. \textbf{158} (2003), no.~2, 635--659.
\href{https://doi.org/10.4007/annals.2003.158.635}{doi:10.4007/annals.2003.158.635}.

\bibitem[CS95]{ChristensenSinclair1995}
E.~Christensen and A.~M.~Sinclair,
\emph{Module mappings into von Neumann algebras and injectivity},
Proc. London Math. Soc. (3) \textbf{71} (1995), no.~3, 618--640.
\href{https://doi.org/10.1112/plms/s3-71.3.618}{doi:10.1112/plms/s3-71.3.618}.

\bibitem[Dix53]{Dixmier1953}
J.~Dixmier,
\emph{Formes lin\'eaires sur un anneau d'op\'erateurs},
Bull. Soc. Math. France \textbf{81} (1953), 9--39.
\href{https://doi.org/10.24033/bsmf.1436}{doi:10.24033/bsmf.1436}.

\bibitem[Fur63]{Furstenberg1963}
H.~Furstenberg,
\emph{Noncommuting random products},
Trans. Amer. Math. Soc. \textbf{108} (1963), no.~3, 377--428.
\href{https://doi.org/10.1090/S0002-9947-1963-0163345-0}{doi:10.1090/S0002-9947-1963-0163345-0}.

\bibitem[Haa79]{Haagerup1979}
U.~Haagerup,
\emph{An example of a non nuclear $C^*$-algebra, which has the metric
approximation property},
Invent. Math. \textbf{50} (1979), 279--293.
\href{https://doi.org/10.1007/BF01410082}{doi:10.1007/BF01410082}.

\bibitem[JKR72]{JKR1972}
B.~E.~Johnson, R.~V.~Kadison, and J.~R.~Ringrose,
\emph{Cohomology of operator algebras. III: Reduction to normal cohomology},
Bull. Soc. Math. France \textbf{100} (1972), 73--96.
\href{https://doi.org/10.24033/bsmf.1731}{doi:10.24033/bsmf.1731}.

\bibitem[Kad66]{Kadison1966}
R.~V.~Kadison,
\emph{Derivations of operator algebras},
Ann. of Math. (2) \textbf{83} (1966), no.~2, 280--293.
\href{https://doi.org/10.2307/1970433}{doi:10.2307/1970433}.

\bibitem[KR71a]{KadisonRingrose1971}
R.~V.~Kadison and J.~R.~Ringrose,
\emph{Cohomology of operator algebras. I: Type I von Neumann algebras},
Acta Math. \textbf{126} (1971), 227--243.
\href{https://doi.org/10.1007/BF02392032}{doi:10.1007/BF02392032}.

\bibitem[KR71b]{KadisonRingrose1971II}
R.~V.~Kadison and J.~R.~Ringrose,
\emph{Cohomology of operator algebras II. Extended cobounding and
the hyperfinite case},
Ark. Mat. \textbf{9} (1971), 55--63.
\href{https://doi.org/10.1007/BF02383637}{doi:10.1007/BF02383637}.

\bibitem[PS10]{PopSmith2010}
F.~Pop and R.~R.~Smith,
\emph{Vanishing of second cohomology for tensor products of type
$\mathrm{II}_1$ von Neumann algebras},
J. Funct. Anal. \textbf{258} (2010), 2695--2707.
\href{https://doi.org/10.1016/j.jfa.2010.01.013}{doi:10.1016/j.jfa.2010.01.013}.

\bibitem[Pop95]{Popa1995}
S.~Popa,
\emph{Free-independent sequences in type $\mathrm{II}_1$ factors
and related problems},
in \emph{Recent advances in operator algebras---Orl\'eans, 1992},
Ast\'erisque \textbf{232} (1995), 187--202.
\href{https://doi.org/10.24033/ast.318}{doi:10.24033/ast.318}.

\bibitem[Pop14]{Popa2014}
S.~Popa,
\emph{Independence properties in subalgebras of ultraproduct
$\mathrm{II}_1$ factors},
J. Funct. Anal. \textbf{266} (2014), no.~9, 5818--5846.
\href{https://doi.org/10.1016/j.jfa.2014.02.004}{doi:10.1016/j.jfa.2014.02.004}.

\bibitem[QS15]{QianShen2015}
W.~Qian and J.~Shen,
\emph{Hochschild cohomology of type $\mathrm{II}_1$ von Neumann
algebras with Property $\Gamma$},
Oper. Matrices \textbf{9} (2015), no.~3, 507--543.
\href{https://doi.org/10.7153/oam-09-31}{doi:10.7153/oam-09-31}.

\bibitem[Sak66]{Sakai1966}
S.~Sakai,
\emph{Derivations of $W^*$-algebras},
Ann. of Math. (2) \textbf{83} (1966), no.~2, 273--279.
\href{https://doi.org/10.2307/1970432}{doi:10.2307/1970432}.

\bibitem[SS95]{SinclairSmith1995}
A.~M.~Sinclair and R.~R.~Smith,
\emph{Hochschild Cohomology of von Neumann Algebras},
London Mathematical Society Lecture Note Series, vol.~203,
Cambridge University Press, Cambridge, 1995.
\href{https://doi.org/10.1017/CBO9780511526190}{doi:10.1017/CBO9780511526190}.

\bibitem[SS98a]{SinclairSmith1998Cartan}
A.~M.~Sinclair and R.~R.~Smith,
\emph{Hochschild cohomology for von Neumann algebras with Cartan subalgebras},
Amer. J. Math. \textbf{120} (1998), no.~5, 1043--1057.
\href{https://doi.org/10.1353/ajm.1998.0042}{doi:10.1353/ajm.1998.0042}.

\bibitem[SS98b]{SinclairSmith1998}
A.~M.~Sinclair and R.~R.~Smith,
\emph{The Hochschild cohomology problem for von Neumann algebras},
Proc. Natl. Acad. Sci. USA \textbf{95} (1998), no.~7, 3376--3379.
\href{https://doi.org/10.1073/pnas.95.7.3376}{doi:10.1073/pnas.95.7.3376}.

\bibitem[Ume54]{Umegaki1954}
H.~Umegaki,
\emph{Conditional expectation in an operator algebra},
Tohoku Math. J. (2) \textbf{6} (1954), nos.~2--3, 177--181.
\href{https://doi.org/10.2748/tmj/1178245177}{doi:10.2748/tmj/1178245177}.

\bibitem[Voi90]{Voiculescu1990Circular}
D.~Voiculescu,
\emph{Circular and semicircular systems and free product factors},
in \emph{Operator Algebras, Unitary Representations, Enveloping
Algebras, and Invariant Theory},
Progress in Mathematics, vol.~92, Birkh\"auser, Boston, 1990, 45--60.
\href{https://math.berkeley.edu/~dvv/CircularSemi.pdf}{Author's text}.

\end{thebibliography}

\end{document}
